% ============================================================
\documentclass[master.tex]{subfiles}

\begin{document}

% ------------------------------------------------------------
%  Title block (standalone) / chapter heading (master)
% ------------------------------------------------------------
\ifSubfilesClassLoaded{%
  \begin{center}
    {\LARGE\bfseries\color{isublue}
     Chapter 12: Estimation under Instrumental Variables}\\[6pt]
    {\large Theory and Methods in Causal Inference}\\[4pt]
    {\normalsize Department of Statistics, Iowa State University}
  \end{center}
  \bigskip
  \hrule height 1.2pt \relax
  \smallskip
}{%
  \chapter{Estimation under Instrumental Variables}
}

% ------------------------------------------------------------
%  Learning objectives
% ------------------------------------------------------------
\begin{tcolorbox}[defstyle, title={Learning Objectives --- Core Material
  (\S\S\,\ref{w12:sec:wald}--\ref{w12:sec:asymp},
   \ref{w12:sec:weak})}]
By the end of the core sections, students should be able to:
\begin{enumerate}
  \item Construct the Wald estimator as the sample analog of the Wald
    estimand and interpret its numerator and denominator as the
    reduced form and first stage.  Extend this to the IV regression
    estimator with covariates as the direct sample analog of
    $\mathrm{Cov}(Y,Z\mid X)/\mathrm{Cov}(T,Z\mid X)$, and recognize
    the Wald estimator as a special case.
  \item Distinguish the structural form of a linear SEM from its
    reduced form; derive the reduced form by solving out endogenous
    variables; interpret the reduced-form coefficient $\phi = \beta\pi$
    as the product of the structural parameter and the first-stage
    strength; define the reduced form regression estimator as OLS
    applied to the reduced form equation; and interpret
    $\hat\phi_{\mathrm{RF}}$ as the intent-to-treat (ITT) effect when
    the instrument is randomly assigned.
  \item Derive the 2SLS estimator from first principles via the
    constrained normal equations argument, and prove its numerical
    equivalence with the IV regression estimator.
  \item Interpret 2SLS as a method-of-moments estimator based on the
    IV orthogonality condition, connect it to the estimating-equation
    framework of Chapter~9, and extend to GMM for overidentified
    models.
  \item State the asymptotic distribution of the GMM estimator,
    specialize the sandwich variance formula to the exactly identified
    and efficient GMM cases, and construct cluster-robust standard
    errors.
  \item Explain the weak-instrument problem and its consequences for
    finite-sample bias and inferential reliability.
\end{enumerate}
\end{tcolorbox}

\begin{tcolorbox}[remarkstyle, title={Learning Objectives --- Advanced
  Enrichment (\S\S\,\ref{w12:sec:gel}--\ref{w12:sec:cf})}]
Students who complete the advanced sections should additionally be
able to:
\begin{enumerate}
  \item Describe the GEL framework as a one-step alternative to
    efficient GMM; state the saddle-point formulation; identify
    empirical likelihood, exponential tilting, and the continuous
    updating estimator as special cases; and interpret the
    Legendre--Fenchel duality as observation re-weighting.
  \item Describe the control function approach, explain its
    equivalence to 2SLS in the linear model, and state the conditions
    under which $V = F_{T \mid Z}(T, Z)$ serves as a control
    variable in the nonlinear triangular model of \citet{imbens2009}.
\end{enumerate}
\end{tcolorbox}


% ============================================================
\section{From Identification to Estimation}
\label{w12:sec:intro}
% ============================================================

Chapter~7 established what IV identifies and under what assumptions.
This chapter shows how to estimate that target from finite data using
sample analogs of the same orthogonality restrictions.

Chapter~7 showed that when an instrument $Z$ satisfies relevance,
exogeneity, and exclusion, the Wald ratio
$\mathrm{Cov}(Y,Z\mid X)/\mathrm{Cov}(T,Z\mid X)$ is a well-defined
function of observable quantities --- but it does not identify any
specific causal parameter without a fourth structural assumption
(Remark~\ref{w7:rem:three-not-enough}).  This chapter works within
the linear constant-effect model
\begin{align}
  Y &= \alpha + \beta T + \gamma^\top X + \varepsilon,
  \label{w12:eq:structural}\\
  T &= \pi Z + \delta^\top X + \eta,
  \label{w12:eq:firststage}
\end{align}
in which the constant-effect restriction ($Y_i(t)-Y_i(0) = \beta$
for all $i$) serves as that fourth assumption.  Under the three core
assumptions plus this restriction, the structural coefficient $\beta$
is identified by the Wald formula:
\[
  \beta
  = \frac{\mathrm{Cov}(Y,\,Z \mid X)}{\mathrm{Cov}(T,\,Z \mid X)}.
\]
Under heterogeneous treatment effects, the same ratio identifies the
LATE for the complier population --- provided the additional
monotonicity assumption ($T_i(1) \geq T_i(0)$ for all $i$) serves
as the fourth structural assumption instead; that interpretive point
is developed in Chapter~7.
In the simplest binary-instrument, no-covariate case, this reduces to
the Wald estimand:
\[
  \beta
  = \frac{\E[Y \mid Z{=}1] - \E[Y \mid Z{=}0]}
         {\E[T \mid Z{=}1] - \E[T \mid Z{=}0]}.
\]

Identification establishes that $\beta$ is a well-defined functional
of the observable distribution $P(Y, T, Z, X)$.  This chapter
addresses the estimation problem: how to recover $\beta$ from a
finite sample.  The answer is not a single formula but a family of
estimators---the Wald estimator, two-stage least squares (2SLS), and
the generalized method of moments (GMM)---each of which is a sample
analog of the same underlying orthogonality restriction, viewed from
a different angle.  This chapter develops estimators for the IV
estimands identified in Chapter~7; it does not weaken or replace the
substantive assumptions needed for identification.

This chapter assumes that a valid IV design has already been
justified.  Our goal here is not to discover instruments, but to
estimate an already-identified IV estimand and quantify its
uncertainty.

The chapter divides into two tiers.

\medskip\noindent\textbf{Core material
(\S\S\,\ref{w12:sec:wald}--\ref{w12:sec:asymp},
 \ref{w12:sec:weak}).}
Sections~\ref{w12:sec:wald}--\ref{w12:sec:asymp} develop the
mainstream IV estimators and their asymptotic theory.
Section~\ref{w12:sec:wald} introduces the Wald estimator and derives
the IV regression estimator with covariates via the constrained normal
equations argument.
Section~\ref{w12:sec:2sls} extends this to two-stage least squares
for multi-valued and multiple instruments.
Section~\ref{w12:sec:equiv} proves the equivalence of the IV
regression estimator and 2SLS.
Section~\ref{w12:sec:gmm} interprets both as method-of-moments
estimators and extends to GMM for overidentified models.
Section~\ref{w12:sec:asymp} develops the asymptotic distribution of
the GMM estimator, specializes the sandwich variance formula, and
covers cluster-robust inference.
Section~\ref{w12:sec:weak} treats the weak-instrument problem.

\medskip\noindent\textbf{Advanced enrichment
(\S\S\,\ref{w12:sec:gel}--\ref{w12:sec:cf}).}
Section~\ref{w12:sec:gel} introduces generalized empirical likelihood
(GEL) as a one-step alternative to efficient GMM, develops the
convex-conjugate duality with the minimum-discrepancy primal problem,
and discusses finite-sample advantages.
Section~\ref{w12:sec:cf} develops the control function approach and
its nonlinear extension due to \citet{imbens2009}.
These sections presuppose the core material but can be omitted
without loss of continuity.


% ============================================================
\section{The Wald Estimator and the IV Regression Estimator}
\label{w12:sec:wald}
% ============================================================

The Wald estimator is the canonical just-identified special case.
It is the cleanest entry point for understanding IV estimation, but
most empirical applications with covariates or multiple instruments
are implemented through regression or GMM formulations.  We begin
with the Wald estimator to establish the sample-analog logic before
generalizing.

\subsection{The Wald Estimator}

We begin with the simplest setting: a binary instrument
$Z \in \{0,1\}$, a scalar treatment $T$, a scalar outcome $Y$, and
no additional covariates $X$.  Chapter~7 established that under
relevance, exogeneity, exclusion, and the constant-effect assumption
($\tau_i = \beta$ for all $i$), the parameter $\beta$ is
identified by the Wald estimand
\[
  \beta
  = \frac{\E[Y \mid Z{=}1] - \E[Y \mid Z{=}0]}
         {\E[T \mid Z{=}1] - \E[T \mid Z{=}0]}.
\]
Each expectation in this ratio is a population mean conditional on
instrument status.  Replacing each with its sample analog yields the
Wald estimator.

\begin{defn}{Wald Estimator}
\label{w12:defn:wald}
Let $n_z = \sum_{i=1}^n \mathbf{1}(Z_i = z)$ for $z \in \{0,1\}$,
and define the within-group sample means
\[
  \bar{Y}_z
  = \frac{1}{n_z}\sum_{i:\,Z_i=z} Y_i,
  \qquad
  \bar{T}_z
  = \frac{1}{n_z}\sum_{i:\,Z_i=z} T_i.
\]
The \emph{Wald estimator} is
\begin{equation}
  \hat\beta_{\mathrm{Wald}}
  = \frac{\bar{Y}_1 - \bar{Y}_0}{\bar{T}_1 - \bar{T}_0}.
  \label{w12:eq:wald-est}
\end{equation}
\end{defn}

Under heterogeneous treatment effects, the Wald estimator remains
statistically valid for the IV estimand, but that estimand is
generally the LATE rather than the ATE.

The Wald estimator has a transparent structural interpretation.
The numerator $\bar{Y}_1 - \bar{Y}_0$ estimates the \emph{reduced
form}: the total effect of the instrument on the outcome.  The
denominator $\bar{T}_1 - \bar{T}_0$ estimates the \emph{first stage}:
the effect of the instrument on treatment uptake.  Their ratio
recovers the effect of treatment on the outcome by attributing all
of the instrument's effect on $Y$ to the path $Z \to T \to Y$---a
step that is valid precisely because the exclusion restriction rules
out any direct path from $Z$ to $Y$.

The three regression objects serve distinct purposes: the first stage
answers how strongly the instrument moves treatment; the reduced form
answers how the instrument changes the outcome; and the Wald (or 2SLS)
ratio converts that intent-to-treat effect into a per-unit-of-treatment
causal effect under the IV assumptions.  All three should be reported
in applied work---not because they are redundant, but because each is
informative about a different aspect of the design.

\begin{remark}[Consistency by the continuous mapping theorem]
The within-group means $\bar{Y}_z$ and $\bar{T}_z$ are consistent
estimators of $\E[Y \mid Z{=}z]$ and $\E[T \mid Z{=}z]$ by the law
of large numbers.  Consistency of $\hat\beta_{\mathrm{Wald}}$ for
$\beta$ then follows from the continuous mapping theorem, provided
the denominator $\E[T \mid Z{=}1] - \E[T \mid Z{=}0] \neq 0$
(the relevance condition).
\end{remark}

\begin{remark}[What the Wald estimator targets under heterogeneous effects]
When treatment effects are heterogeneous, Chapter~7 showed that under
the additional monotonicity assumption ($T_i(1) \geq T_i(0)$ for all
$i$, ruling out defiers) --- the fourth structural assumption in
Framework~2 --- the Wald estimand identifies the Local Average
Treatment Effect (LATE) for the subpopulation of \emph{compliers}:
units whose treatment status is changed by the instrument.  The Wald
estimator therefore consistently estimates the LATE, not the population
average treatment effect.  This interpretive point is separate from,
and does not affect, the statistical validity of the estimator.
\end{remark}

\subsection{The IV Regression Estimator with Covariates}

When observed covariates $X \in \mathbb{R}^p$ are present, the Wald
estimator is no longer applicable: the binary group-mean difference
does not partial out $X$.  We derive the general IV estimator
directly from the estimating-equation principle: an exogenous
variable is valid as an instrument precisely because it satisfies an
orthogonality condition with the structural error that the endogenous
treatment does not.

\medskip\noindent\textbf{Notation.}
Stack observations into matrices $\mathbf{Y}, \mathbf{T}
\in \mathbb{R}^n$, $\mathbf{X} \in \mathbb{R}^{n \times p}$,
$\mathbf{Z} \in \mathbb{R}^n$ (scalar instrument).  Define the
\emph{annihilator matrix}
\begin{equation}
  M_X = I_n - \mathbf{X}(\mathbf{X}^\top\mathbf{X})^{-1}\mathbf{X}^\top,
  \label{w12:eq:mx}
\end{equation}
the orthogonal projection onto the orthogonal complement of
$\mathrm{col}(\mathbf{X})$.  For any vector $\mathbf{W} \in
\mathbb{R}^n$, $M_X\mathbf{W}$ is the vector of OLS residuals from
projecting $\mathbf{W}$ on $\mathbf{X}$; equivalently, the $i$-th
entry of $M_X\mathbf{W}$ is the within-$X$ residual $\tilde{W}_i$.
The matrix $M_X$ is symmetric and idempotent: $M_X^\top = M_X$ and
$M_X^2 = M_X$.

\medskip\noindent\textbf{Derivation via constrained normal equations.}
Consider the augmented model
\begin{equation}
  \mathbf{Y}
  = \mathbf{T}\beta + \mathbf{X}\gamma + \mathbf{Z}\phi + \mathbf{u},
  \qquad \E[\mathbf{u} \mid \mathbf{T}, \mathbf{X}, \mathbf{Z}]
  = \mathbf{0},
  \label{w12:eq:augmented}
\end{equation}
subject to the exclusion restriction $\phi = 0$.  If we were to run
OLS on~\eqref{w12:eq:augmented}, the normal equations for the three
blocks of regressors, with residual $\hat{\mathbf{e}} = \mathbf{Y} -
\mathbf{T}\beta - \mathbf{X}\gamma$, would be:
\begin{align}
  \mathbf{T}^\top \hat{\mathbf{e}} &= \mathbf{0}, \tag{NE-T}\\
  \mathbf{X}^\top \hat{\mathbf{e}} &= \mathbf{0}, \tag{NE-X}\\
  \mathbf{Z}^\top \hat{\mathbf{e}} &= \mathbf{0}. \tag{NE-Z}
\end{align}
At the true $\beta$, the population version of (NE-T) fails:
$\E[\mathbf{T}^\top\boldsymbol\varepsilon] \neq 0$ because $T$ is
endogenous.  By contrast, (NE-X) is valid ($\E[\mathbf{X}^\top
\boldsymbol\varepsilon] = \mathbf{0}$ by exogeneity of $X$) and
(NE-Z) is valid ($\E[\mathbf{Z}^\top\boldsymbol\varepsilon] =
\mathbf{0}$ by exogeneity and exclusion of $Z$).

We therefore \emph{drop} the invalid (NE-T) and solve the
$(p+1)$-dimensional system (NE-X), (NE-Z).  From (NE-X):
\[
  \hat\gamma
  = (\mathbf{X}^\top\mathbf{X})^{-1}
    (\mathbf{X}^\top\mathbf{Y} - \mathbf{X}^\top\mathbf{T}\beta).
\]
Substituting into (NE-Z) and simplifying using $M_X$:
\[
  \mathbf{Z}^\top M_X \mathbf{T}\;\beta
  = \mathbf{Z}^\top M_X \mathbf{Y}.
\]
Provided $\mathbf{Z}^\top M_X \mathbf{T} \neq 0$ (the relevance
condition after partialling out $X$), this has the unique solution

\begin{defn}{IV Regression Estimator}
\label{w12:defn:iv-reg}
The \emph{IV regression estimator} in the linear IV model with
scalar instrument $Z$ and covariates $X$ is
\begin{equation}
  \hat\beta_{\mathrm{IV}}
  = \frac{\mathbf{Z}^\top M_X \mathbf{Y}}
         {\mathbf{Z}^\top M_X \mathbf{T}}
  = \frac{\widehat{\mathrm{Cov}}(\tilde{Z},\, Y)}
         {\widehat{\mathrm{Cov}}(\tilde{Z},\, T)},
  \label{w12:eq:iv-reg}
\end{equation}
where $\tilde{Z}_i$, $\tilde{Y}_i$, $\tilde{T}_i$ are the within-$X$
residuals (entries of $M_X\mathbf{Z}$, $M_X\mathbf{Y}$,
$M_X\mathbf{T}$).
\end{defn}

The estimator forces the residual $\hat{\mathbf{e}} = \mathbf{Y} -
\mathbf{T}\hat\beta_{\mathrm{IV}} - \mathbf{X}\hat\gamma$ to satisfy
both (NE-X) and (NE-Z); it is the unique solution obtained by
substituting the instrument equation for the corrupted treatment
equation.

\begin{remark}[Wald estimator as a special case]
When $Z \in \{0,1\}$ is binary and there are no covariates ($X$
absent), $M_X = I_n - n^{-1}\mathbf{1}\mathbf{1}^\top$ is the
centering matrix and $M_X\mathbf{W}$ is the mean-centered $\mathbf{W}$.
One can verify that $\mathbf{Z}^\top M_X\mathbf{Y} \propto
(\bar{Y}_1 - \bar{Y}_0)$ and $\mathbf{Z}^\top M_X\mathbf{T} \propto
(\bar{T}_1 - \bar{T}_0)$, so $\hat\beta_{\mathrm{IV}} =
\hat\beta_{\mathrm{Wald}}$.
\end{remark}

\subsection{Structural Form, Reduced Form, and the Reduced Form Regression}
\label{w12:subsec:reduced-form}

The terminology \emph{reduced form} only makes sense in contrast to
something that has not been reduced.  This subsection introduces that
contrast, derives the reduced form of the linear SEM, and then
presents the reduced form regression as the procedure for estimating
the relevant coefficient.

\medskip\noindent\textbf{The structural form.}
The system~\eqref{w12:eq:structural}--\eqref{w12:eq:firststage} is
called the \emph{structural form} of the model.  It represents the
causal mechanism directly: each equation describes how one variable
is determined by others, including endogenous variables on the
right-hand side.  In particular, the outcome equation
\[
  Y = \alpha + \beta T + \gamma^\top X + \varepsilon
\]
contains the endogenous regressor $T$ on its right-hand side.
The structural coefficient $\beta$ has a causal interpretation
(it is the effect of a unit change in $T$ on $Y$, holding $X$ and
$\varepsilon$ fixed), but $T$ is correlated with $\varepsilon$,
so OLS applied to the structural form is inconsistent.

\medskip\noindent\textbf{The reduced form.}
The \emph{reduced form} is obtained by solving the structural system
so that each endogenous variable is expressed purely as a function
of the exogenous variables $Z$ and $X$, with no endogenous
regressors on the right-hand side.  Substituting the first-stage
equation~\eqref{w12:eq:firststage} into the outcome
equation~\eqref{w12:eq:structural} gives
\begin{equation}
  Y = \underbrace{(\alpha + \beta\delta^\top X)}_{\text{intercept/}X\text{-terms}}
    + \underbrace{\beta\pi}_{\phi}\, Z
    + (\beta\delta + \gamma)^\top X
    + \underbrace{(\beta\eta + \varepsilon)}_{\nu}.
  \label{w12:eq:reduced-form-pop}
\end{equation}
Writing this compactly:
\begin{equation}
  Y = \alpha_{\mathrm{rf}} + \phi\, Z
    + \gamma_{\mathrm{rf}}^\top X + \nu,
  \label{w12:eq:reduced-form}
\end{equation}
where $\phi = \beta\pi$, $\gamma_{\mathrm{rf}} = \beta\delta + \gamma$,
and $\nu = \beta\eta + \varepsilon$.  The reduced-form equation for
$T$ is simply the first-stage equation~\eqref{w12:eq:firststage}
itself, which is already in reduced form because $Z$ and $X$ are
exogenous.

Two features of the reduced form are immediate from
\eqref{w12:eq:reduced-form-pop}.  First, both right-hand-side
variables ($Z$ and $X$) are exogenous, so OLS applied
to~\eqref{w12:eq:reduced-form} is consistent for $\phi$.  Second,
$\phi = \beta\pi$ is a product of the structural parameter of
interest $\beta$ and the first-stage strength $\pi$.  Neither
$\beta$ nor $\pi$ is separately identified from the reduced form
alone; identifying $\beta$ requires also using the first-stage
equation---which is the essence of IV.

Because the exclusion restriction rules out any direct path
$Z \to Y$, the full effect of the instrument on the outcome flows
through the treatment channel $Z \to T \to Y$, so $\phi = \beta\pi$
and not $\phi = \beta\pi + \delta$ for some additional direct
effect $\delta$.

\medskip\noindent\textbf{The reduced form regression.}
The reduced form regression is simply OLS applied to the reduced
form equation~\eqref{w12:eq:reduced-form}.  Its purpose is to
estimate the reduced-form coefficient $\phi$.

\begin{defn}{Reduced Form Regression Estimator}
\label{w12:defn:reduced-form}
The \textbf{reduced form regression estimator} is the OLS estimator
of $\phi$ in the reduced form equation~\eqref{w12:eq:reduced-form}:
\begin{equation}
  \hat\phi_{\mathrm{RF}}
  = \frac{\widehat{\mathrm{Cov}}(\tilde{Z},\, Y)}
         {\widehat{\mathrm{Var}}(\tilde{Z})},
  \label{w12:eq:rf-est}
\end{equation}
where $\tilde{Z}$ is the residual from projecting $Z$ on
$(1, X^\top)^\top$.  It is a consistent estimator of
$\phi = \beta\pi$ under the IV exogeneity condition
$\E[\nu \mid Z, X] = 0$, which holds whenever exogeneity and
exclusion hold for the structural form.
\end{defn}

\medskip\noindent\textbf{Intent-to-treat interpretation.}
When $Z$ is randomly assigned (as in a lottery or randomized
encouragement design), $\hat\phi_{\mathrm{RF}}$ estimates the
\emph{intent-to-treat} (ITT) effect: the average effect on $Y$ of
being assigned $Z = 1$ rather than $Z = 0$, regardless of whether
treatment is actually taken up.  The ITT is a valid causal estimand
in its own right---it requires only exogeneity of $Z$, not the
exclusion restriction---and it is often of direct policy interest
when the instrument corresponds to an implementable intervention.

\medskip\noindent\textbf{Relationship to the first stage and the IV
estimator.}
The first-stage regression estimator is
\[
  \hat\pi_{\mathrm{FS}}
  = \frac{\widehat{\mathrm{Cov}}(\tilde{Z},\, T)}
         {\widehat{\mathrm{Var}}(\tilde{Z})}.
\]
Since $\phi = \beta\pi$, we have $\beta = \phi/\pi$, and the sample
analog gives
\begin{equation}
  \hat\beta_{\mathrm{IV}}
  = \frac{\hat\phi_{\mathrm{RF}}}{\hat\pi_{\mathrm{FS}}},
  \label{w12:eq:iv-ratio}
\end{equation}
which reproduces the IV regression estimator~\eqref{w12:eq:iv-reg}.
The reduced form delivers the instrument's total effect on the
outcome; the first stage scales it by the instrument's effect on
treatment; the ratio removes the scaling and recovers the structural
parameter.  Equation~\eqref{w12:eq:iv-ratio} also shows that the IV
estimator is undefined when $\hat\pi_{\mathrm{FS}} = 0$---the
sample-level manifestation of the relevance requirement.

The ratio form~\eqref{w12:eq:iv-ratio} is useful because it exposes
the logic of IV, not because all IV estimation should be carried out
by separately estimating and then dividing two coefficients.  In
practice, 2SLS (Section~\ref{w12:sec:2sls}) solves the estimating
equations jointly and produces valid standard errors automatically;
forming $\hat\phi_{\mathrm{RF}}/\hat\pi_{\mathrm{FS}}$ by hand and
dividing their separate standard errors is incorrect.

\begin{remark}[Reporting all three quantities]
In applied work it is standard practice to report the first stage,
reduced form, and IV (or 2SLS) estimates side by side, as in
Table~7.3 of Chapter~7.  The reduced form and first stage each have
a transparent OLS interpretation and can be assessed independently
before the ratio is formed.  Reporting all three allows the reader
to verify instrument strength (first stage), the ITT effect (reduced
form), and the per-unit-of-treatment causal effect (IV/2SLS)
simultaneously.  The IV estimate carries no identifying content
beyond what the reduced form and first stage together contain.
\end{remark}


% ============================================================
\section{Two-Stage Least Squares}
\label{w12:sec:2sls}
% ============================================================

The Wald estimator is limited to binary instruments and no covariates.
Two-stage least squares (2SLS) extends the identification-to-estimation
step to settings with continuous or multi-valued instruments, multiple
instruments, and observed covariates.  The key idea is that 2SLS
replaces the endogenous regressor by its projection onto the instrument
space, thereby retaining only the variation in treatment that is
explained by the exogenous instruments and covariates.  Although it
is defined through a two-stage procedure below, 2SLS is a single
estimator with a closed-form expression; the stages are a transparent
derivation device, not a description of two separate estimation
problems.

\subsection{Setup}

Return to the structural model~\eqref{w12:eq:structural}--\eqref{w12:eq:firststage},
where $T$ is a scalar endogenous variable, $X \in \mathbb{R}^p$ is a
vector of exogenous covariates, and $Z \in \mathbb{R}^q$ is a vector of
instruments with $q \geq 1$.  The IV assumptions require:
\begin{enumerate}
  \item \textbf{Relevance}: the first-stage coefficient $\pi \neq 0$,
    i.e., $\mathrm{Cov}(Z, T \mid X) \neq 0$.
  \item \textbf{Exogeneity and exclusion}: $\E[\varepsilon \mid Z, X] = 0$.
\end{enumerate}
Endogeneity means $\mathrm{Cov}(T, \varepsilon) \neq 0$, arising
because $\varepsilon$ and $\eta$ share an unobserved common driver.

\subsection{Stage 1: The First-Stage Regression}

Regress $T$ on $Z$ and $X$ by OLS:
\[
  T = \pi^\top Z + \delta^\top X + \eta.
\]
Denote the OLS fitted values
\[
  \hat{T}_i = \hat\pi^\top Z_i + \hat\delta^\top X_i.
\]
The fitted value $\hat{T}_i$ is the linear projection of $T_i$ onto
the space spanned by $Z_i$ and $X_i$.  Because OLS residuals are
orthogonal to the regressors, $\hat{T}_i$ is by construction
uncorrelated (in the sample) with the first-stage residual
$\hat\eta_i = T_i - \hat{T}_i$.  The fitted values $\hat{T}_i$ are
not used because they predict treatment well in a generic regression
sense; they are used because they isolate the component of treatment
variation spanned by the instruments and covariates---the variation
that is exogenous under IV validity.

\subsection{Stage 2: The Second-Stage Regression}

Regress $Y$ on $\hat{T}$ and $X$ by OLS:
\[
  Y = \alpha + \beta \hat{T} + \gamma^\top X + \text{error}.
\]

\begin{defn}{2SLS Estimator}
\label{w12:defn:2sls}
The \emph{two-stage least squares estimator} $\hat\beta_{\mathrm{2SLS}}$
is the OLS coefficient on $\hat{T}$ in the second-stage regression.
\end{defn}

Equivalently, 2SLS projects the endogenous regressors onto the linear
span of the instruments and exogenous covariates, and then uses only
that projected component in the structural regression.  The
endogenous part of $T$---the residual $\hat\eta_i = T_i -
\hat{T}_i$ that is correlated with $\varepsilon_i$---is discarded
before the causal coefficient is estimated.

To obtain a closed form, let $\tilde{Y}_i$, $\tilde{T}_i$, and
$\tilde{Z}_i$ denote the residuals from projecting $Y_i$, $T_i$, and
$Z_i$, respectively, onto $X_i$ by OLS (i.e., the within-$X$-residuals).
In the single-instrument case ($q = 1$), the 2SLS estimator simplifies to
\begin{equation}
  \hat\beta_{\mathrm{2SLS}}
  = \frac{\sum_{i=1}^n \tilde{Z}_i\, Y_i}
         {\sum_{i=1}^n \tilde{Z}_i\, T_i}
  = \frac{\widehat{\mathrm{Cov}}(\tilde{Z},\, Y)}
         {\widehat{\mathrm{Cov}}(\tilde{Z},\, T)}.
  \label{w12:eq:2sls-closed}
\end{equation}
This is structurally parallel to the Wald ratio: a ratio of
the instrument-by-outcome sample covariance to the
instrument-by-treatment sample covariance, after partialling out the
observed covariates $X$.

\begin{warn}{Standard Errors from the Second-Stage OLS Are Incorrect}
A common error is to report the standard errors from the second-stage
OLS regression directly.  Those standard errors use $\hat{T}$ rather
than $T$ and residuals that do not equal the structural errors
$\varepsilon_i$, so the resulting standard errors are invalid.
Correct inference requires the sandwich variance formula derived in
Section~\ref{w12:sec:asymp}, or software that implements 2SLS natively.
\end{warn}


% ============================================================
\section{Equivalence of the IV Regression Estimator and 2SLS}
\label{w12:sec:equiv}
% ============================================================

The IV regression estimator~\eqref{w12:eq:iv-reg} and the 2SLS
estimator look like different procedures: one computes a ratio of
partialled-out covariances directly, while the other constructs
first-stage fitted values and runs a second regression.  In the
single-instrument case they are in fact numerically identical for any
covariate vector $X$.  This theorem is not just algebraic
housekeeping: it shows that the ratio-of-covariances view and the
staged-regression view are two computational faces of the same IV
estimator.  Knowing both is useful---the ratio form exposes the logic
of IV estimation while the staged-regression form generalizes to
multiple instruments and connects naturally to GMM.

\begin{thm}{IV Regression--2SLS Equivalence}
\label{w12:thm:equiv}
In the linear IV model with a scalar instrument $Z$ and covariates
$X \in \mathbb{R}^p$, the IV regression estimator and the 2SLS
estimator coincide:
\[
  \hat\beta_{\mathrm{2SLS}} = \hat\beta_{\mathrm{IV}}.
\]
\end{thm}

The intuition is compact.  In §\ref{w12:sec:wald} we derived
$\hat\beta_{\mathrm{IV}}$ by dropping the invalid (NE-T) and solving
the valid system (NE-X), (NE-Z), obtaining
$\hat\beta_{\mathrm{IV}} = (\mathbf{Z}^\top M_X \mathbf{T})^{-1}
\mathbf{Z}^\top M_X \mathbf{Y}$.
The 2SLS procedure reaches the same formula by a projection
argument.

\begin{proof}
Retain the matrix notation of \S\ref{w12:sec:wald}: $M_X =
I_n - \mathbf{X}(\mathbf{X}^\top\mathbf{X})^{-1}\mathbf{X}^\top$,
scalar instrument $\mathbf{Z} \in \mathbb{R}^n$.

\textbf{First-stage projection.}
The first-stage fitted value is $\hat{\mathbf{T}} =
P_{[\mathbf{X},\mathbf{Z}]}\mathbf{T}$, where
$P_{[\mathbf{X},\mathbf{Z}]}$ projects onto
$\mathrm{col}(\mathbf{X}, \mathbf{Z})$.  By the
Frisch--Waugh--Lovell theorem applied to Stage~1:
\[
  M_X\hat{\mathbf{T}}
  = M_X\mathbf{Z}
    (\mathbf{Z}^\top M_X\mathbf{Z})^{-1}
    \mathbf{Z}^\top M_X\mathbf{T}.
\]

\textbf{2SLS second-stage closed form.}
Applying FWL to the second-stage regression of $\mathbf{Y}$ on
$(\hat{\mathbf{T}}, \mathbf{X})$:
\[
  \hat\beta_{\mathrm{2SLS}}
  = \frac{\hat{\mathbf{T}}^\top M_X \mathbf{Y}}
         {\hat{\mathbf{T}}^\top M_X \hat{\mathbf{T}}}
  = \frac{\mathbf{T}^\top M_X\mathbf{Z}
          (\mathbf{Z}^\top M_X\mathbf{Z})^{-1}
          \mathbf{Z}^\top M_X\mathbf{Y}}
         {\mathbf{T}^\top M_X\mathbf{Z}
          (\mathbf{Z}^\top M_X\mathbf{Z})^{-1}
          \mathbf{Z}^\top M_X\mathbf{T}}.
\]
Under exact identification ($q = 1$), the scalars
$\mathbf{T}^\top M_X\mathbf{Z}$ and
$(\mathbf{Z}^\top M_X\mathbf{Z})^{-1}$ cancel from numerator and
denominator, leaving
\[
  \hat\beta_{\mathrm{2SLS}}
  = \frac{\mathbf{Z}^\top M_X\mathbf{Y}}
         {\mathbf{Z}^\top M_X\mathbf{T}}
  = \hat\beta_{\mathrm{IV}}. \qedhere
\]
\end{proof}

In the scalar-instrument linear model, Wald, IV regression, and 2SLS
are therefore not competing methods; they are different representations
of the same sample analog of the identification formula.

\begin{remark}[What the proof reveals]
The proof makes the logic of IV transparent at the level of estimating
equations.  In OLS, the normal equation for each regressor forces the
residual to be orthogonal to that regressor.  When $T$ is endogenous,
its normal equation is corrupted: $\E[\mathbf{T}^\top\boldsymbol\varepsilon]
\neq 0$, so the OLS condition $\mathbf{T}^\top\hat{\mathbf{e}} = 0$
holds at the wrong $\theta$.  IV simply replaces this corrupted equation
with the instrument equation $\mathbf{Z}^\top\hat{\mathbf{e}} = 0$,
which is valid because $Z$ is exogenous.  The IV estimator is
therefore OLS on the structural model but with one invalid normal
equation swapped for a valid one.  The 2SLS derivation reaches the
same formula by a projection argument, confirming that the two
approaches are two paths to the same system.
\end{remark}

\begin{remark}[Wald estimator as a special case]
\label{w12:rem:wald-special}
When $Z \in \{0,1\}$ is binary and $X$ is absent, the IV regression
estimator reduces to the Wald estimator
(Definition~\ref{w12:defn:wald}), as shown in
Section~\ref{w12:sec:wald}.  Theorem~\ref{w12:thm:equiv} therefore
implies $\hat\beta_{\mathrm{2SLS}} = \hat\beta_{\mathrm{Wald}}$ in
this case as well.
\end{remark}

\begin{remark}[Generalization to multiple instruments]
With $q > 1$ instruments, the IV regression formula no longer applies
directly because the system $\sum_i \tilde{Z}_i Y_i / \sum_i
\tilde{Z}_i T_i$ is not well defined for vector $\tilde{Z}_i$.  2SLS
resolves this by using the first-stage projection of $T$ onto the
full instrument vector, producing a single scalar $\hat{T}_i$.  The
resulting estimator is efficient in the class of linear IV estimators
under homoskedasticity and reduces to the IV regression estimator when
$q = 1$.
\end{remark}


% ============================================================
\section{The Moment-Condition View and GMM}
\label{w12:sec:gmm}
% ============================================================

The proof of Theorem~\ref{w12:thm:equiv} already contains the key
idea: IV estimation amounts to solving the normal equations of the
structural model with the invalid $T$-block replaced by the valid
instrument equation $\mathbf{Z}^\top\hat{\mathbf{e}} = 0$.  This
replacement is not an algebraic trick but a direct expression of the
IV assumption $\E[Z\varepsilon \mid X] = 0$, which is a
\emph{moment condition}.  The advantage of the moment-condition view
is that it separates the identifying restrictions from the specific
algebra of 2SLS, making clear which part of the problem is about
causal assumptions and which part is about estimation efficiency.
Recognizing 2SLS as a method-of-moments estimator for this condition
connects it to the general estimating-equation framework of Chapter~9
and motivates the generalized method of moments for overidentified
models.

\subsection{2SLS as a Method-of-Moments Estimator}

Write the IV moment condition in terms of the full instrument and
covariate vector $W = (Z^\top, X^\top)^\top \in \mathbb{R}^{q+p}$:
\[
  \E\bigl[W\,\bigl(Y - \alpha - \beta T - \gamma^\top X\bigr)\bigr] = 0.
\]
Setting $\theta = (\alpha, \beta, \gamma^\top)^\top$ and defining the
moment function
\[
  U(O;\,\theta) = W\,(Y - \alpha - \beta T - \gamma^\top X),
\]
the identifying condition is $\E\{U(O;\,\theta_0)\} = 0$, which is
exactly an estimating equation in the sense of Chapter~9.  When the
number of moment conditions equals the number of parameters (exactly
identified: $q + p = 1 + p$, i.e., $q = 1$), the sample-moment
equations $\mathbb{P}_n U(O;\,\hat\theta) = 0$ admit a unique solution.

\begin{thm}{2SLS as a Method-of-Moments Estimator}
\label{w12:thm:mm}
In the exactly identified linear IV model ($q = 1$), the 2SLS estimator
$\hat\beta_{\mathrm{2SLS}}$ is the unique solution to the sample
moment condition
\[
  \frac{1}{n}\sum_{i=1}^n
  Z_i\,\bigl(Y_i - \hat\alpha - \hat\beta_{\mathrm{2SLS}}\,T_i
             - \hat\gamma^\top X_i\bigr)
  = 0,
\]
solved jointly for $(\hat\alpha, \hat\beta_{\mathrm{2SLS}}, \hat\gamma)$.
\end{thm}

This moment-condition view clarifies why the first-stage fitted values
$\hat{T}_i$ are not merely a computational device.  They encode the
orthogonality structure implied by instrument validity: replacing $T_i$
with $\hat{T}_i$ ensures the second-stage regression uses only the
exogenous variation in treatment.  IV is therefore not a separate
estimation philosophy from the rest of Part~III; it is a special case
of the general estimating-equation framework of Chapter~9, with moment
functions induced by instrument validity.

\subsection{Overidentification and GMM}

When $q > 1$ instruments are available, the model is
\emph{overidentified}: there are more moment conditions than
parameters.  The system $\mathbb{P}_n U(O;\,\theta) = 0$ is
generically overdetermined and has no exact solution.

When the model is overidentified, different linear combinations of
the sample moments contain the same target parameter but with
different noise levels.  GMM chooses the combination that solves
the moment equations as well as possible, with a weighting scheme
that can exploit the heteroskedastic structure of the moment scores
for greater efficiency.  The weight matrix $\hat{W}_n$ determines
which combination is used: identity weighting treats all moment
conditions equally, while efficient weighting down-weights noisier
moments proportionally to their variance.

The \emph{generalized method of moments} (GMM) formalizes this by
minimizing a weighted quadratic form in the sample moments:
\[
  \hat\theta_{\mathrm{GMM}}
  = \arg\min_{\theta}\;
  \bigl[\mathbb{P}_n U(O;\,\theta)\bigr]^\top\,
  \hat{W}_n\,
  \bigl[\mathbb{P}_n U(O;\,\theta)\bigr],
\]
where $\hat{W}_n$ is a positive-definite weighting matrix.  Different
choices of $\hat{W}_n$ yield different estimators within the IV class:
\begin{itemize}
  \item A particular choice of $\hat{W}_n$ involving the first-stage
    projection matrix yields the 2SLS estimator.
  \item The choice $\hat{W}_n = \bigl[\mathbb{P}_n U(O;\,\hat\theta)\,
    U(O;\,\hat\theta)^\top\bigr]^{-1}$---the inverse of the estimated
    moment covariance---yields the \emph{efficient GMM estimator},
    which achieves the semiparametric efficiency bound for the
    IV moment-condition model.
\end{itemize}
Under homoskedasticity, efficient GMM and 2SLS coincide.  Under
heteroskedasticity, efficient GMM is strictly more efficient.

\begin{remark}[Overidentification as a testable restriction]
An overidentified model imposes more moment conditions than are
needed for point identification.  If all instruments are valid, the
sample moments should be jointly near zero.  The Sargan--Hansen
$J$-statistic \citep{sargan1958estimation, hansen1982large} tests
this restriction: under the null that all moment conditions hold,
$J \xrightarrow{d} \chi^2_{q-1}$.  A rejection indicates that the
full set of maintained moment conditions is incompatible with the
data.  This may reflect an exclusion failure, an exogeneity failure,
or other model misspecification; the test does not identify which
assumption failed or which instrument is responsible.
Overidentification therefore allows us to check whether the full set
of moment restrictions is jointly compatible with the data, but a
failure to reject is consistent with all instruments being valid, or
with violations that happen to cancel in the joint test.
\end{remark}

\begin{example}{GMM with Two Instruments}
\label{w12:ex:gmm-toy}
We illustrate the GMM procedure in the simplest overidentified
setting: one endogenous regressor, two instruments, no additional
covariates.

\textbf{Setup.}
Consider the structural model $Y = \beta T + \varepsilon$ (intercept
suppressed after demeaning).  Two instruments $Z_1$ and $Z_2$ are
available, giving the moment conditions
\[
  \E[Z_1\varepsilon] = 0,
  \qquad
  \E[Z_2\varepsilon] = 0.
\]
The model is overidentified ($q = 2$, one free parameter), so
the two moment conditions cannot generally be satisfied
simultaneously in finite samples.

\textbf{Sample moments.}
Suppose the following sample covariances are observed:
\[
  \widehat{\mathrm{Cov}}(Z_1,\,Y) = 0.40,
  \quad
  \widehat{\mathrm{Cov}}(Z_1,\,T) = 0.50,
  \quad
  \widehat{\mathrm{Cov}}(Z_2,\,Y) = 0.60,
  \quad
  \widehat{\mathrm{Cov}}(Z_2,\,T) = 0.80.
\]
The just-identified IV estimate using each instrument alone is
\[
  \hat\beta^{(1)} = \frac{0.40}{0.50} = 0.80,
  \qquad
  \hat\beta^{(2)} = \frac{0.60}{0.80} = 0.75.
\]
The two instruments give different point estimates, reflecting
sampling variation and possibly differing local estimands.  The
GMM objective combines both moment conditions.

\textbf{GMM moment function.}
With $a = (\widehat{\mathrm{Cov}}(Z_1,T),\;
\widehat{\mathrm{Cov}}(Z_2,T))^\top = (0.50,\;0.80)^\top$ and
$c = (\widehat{\mathrm{Cov}}(Z_1,Y),\;
\widehat{\mathrm{Cov}}(Z_2,Y))^\top = (0.40,\;0.60)^\top$,
the sample moment vector is
\[
  \hat{U}(\beta)
  = c - a\beta
  = \begin{pmatrix} 0.40 - 0.50\beta \\ 0.60 - 0.80\beta \end{pmatrix}.
\]

\textbf{GMM with identity weighting ($\Omega = I_2$).}
The GMM criterion is $Q(\beta) = \hat{U}(\beta)^\top \hat{U}(\beta)$.
The first-order condition $\partial Q/\partial\beta = 0$ gives
\[
  a^\top(c - a\beta) = 0
  \;\;\Longrightarrow\;\;
  \hat\beta_{\mathrm{GMM}}
  = \frac{a^\top c}{a^\top a}
  = \frac{(0.50)(0.40)+(0.80)(0.60)}{(0.50)^2+(0.80)^2}
  = \frac{0.68}{0.89}
  \approx 0.764.
\]
This is a precision-weighted average of the two just-identified
estimates: it lies between $0.75$ and $0.80$, with more weight on
$Z_2$ because it has the larger first-stage covariance ($0.80 >
0.50$).

\textbf{Effect of the weighting matrix.}
With a diagonal weight $\Omega = \mathrm{diag}(w_1, w_2)$, the
GMM estimator becomes
\[
  \hat\beta_{\mathrm{GMM}}(\Omega)
  = \frac{w_1(0.50)(0.40) + w_2(0.80)(0.60)}
         {w_1(0.50)^2      + w_2(0.80)^2}.
\]
Setting $w_1 = 4$ and $w_2 = 1$ (upweighting $Z_1$) gives
\[
  \hat\beta_{\mathrm{GMM}}
  = \frac{4(0.20)+1(0.48)}{4(0.25)+1(0.64)}
  = \frac{1.28}{1.64}
  \approx 0.780,
\]
closer to $\hat\beta^{(1)} = 0.80$ as expected.  As
$w_1/w_2 \to \infty$, the estimate converges to $\hat\beta^{(1)}$;
as $w_1/w_2 \to 0$, it converges to $\hat\beta^{(2)}$.  The
efficient weighting $\Omega^\ast = \hat\Sigma^{-1}$ (Section
\ref{w12:sec:asymp}) selects $w_1$ and $w_2$ to minimize the
asymptotic variance, typically downweighting the moment condition
with higher residual variance.

\textbf{Sargan--Hansen $J$-test.}
At $\hat\beta \approx 0.764$ under identity weighting, the
residual sample moments are
\[
  \hat{U}_1 = 0.40 - (0.50)(0.764) = 0.018,
  \qquad
  \hat{U}_2 = 0.60 - (0.80)(0.764) = -0.011.
\]
Neither is exactly zero, reflecting the overidentification.  The
$J$-statistic $J = n\,\hat{U}(\hat\beta)^\top\hat\Sigma^{-1}\hat{U}(\hat\beta)$
has a $\chi^2_1$ null distribution (one overidentifying restriction).
A failure to reject is consistent with both instruments satisfying
the exclusion restriction; a rejection points to at least one
invalid instrument.
\end{example}


% ============================================================
\section{Asymptotic Theory of the GMM Estimator}
\label{w12:sec:asymp}
% ============================================================

Section~\ref{w12:sec:gmm} placed IV estimation within the GMM
framework.  This section derives the asymptotic distribution of the
GMM estimator for the linear IV model and then specializes to 2SLS
and to the efficient GMM estimator.

\subsection{Setup and Notation}

Stack the structural regressors and instruments into vectors:
\[
  D_i = \begin{pmatrix} 1 \\ T_i \\ X_i \end{pmatrix} \in \mathbb{R}^k,
  \qquad
  W_i = \begin{pmatrix} 1 \\ X_i \\ Z_i \end{pmatrix} \in \mathbb{R}^m,
\]
where $k = p + 2$ and $m = p + 1 + q$.  The structural model is
$Y_i = D_i^\top\theta_0 + \varepsilon_i$ with
$\theta_0 = (\alpha_0, \beta_0, \gamma_0^\top)^\top$, and the IV
moment condition is
\[
  \E\bigl[U(O_i;\,\theta_0)\bigr] = 0,
  \qquad
  U(O_i;\,\theta) = W_i\,(Y_i - D_i^\top\theta).
\]
The order condition requires $m \geq k$ (at least as many instruments
as parameters), i.e., $q \geq 1$.  For identification, the rank
condition requires $A = \E[W_i D_i^\top] \in \mathbb{R}^{m \times k}$
to have full column rank $k$.

Define the moment variance matrix
\[
  \Sigma = \E\bigl[U(O_i;\,\theta_0)\,U(O_i;\,\theta_0)^\top\bigr]
         = \E[\varepsilon_i^2\, W_i W_i^\top]
         \;\in \mathbb{R}^{m \times m}.
\]
Under homoskedasticity $\E[\varepsilon_i^2 \mid W_i] = \sigma^2$,
this simplifies to $\Sigma = \sigma^2\,\E[W_i W_i^\top]$.

\subsection{Asymptotic Distribution of the GMM Estimator}

The large-sample logic is the same as for any method-of-moments
estimator: sample orthogonality converges to population orthogonality
by the law of large numbers, and the estimator fluctuates around the
truth according to how sensitive the moment equations are to the
parameter and how noisy the sample moments are.  The sensitivity is
captured by the matrix $A = -\E[\partial U/\partial\theta^\top]$,
which measures how sharply the moments respond to a change in
$\theta$; the noise is captured by $\Sigma = \E[U(O;\theta_0)
U(O;\theta_0)^\top]$, the variance of the moment scores at the truth.
The asymptotic variance is a sandwich that combines both.

Recall that the GMM estimator minimizes
$[\mathbb{P}_n U(O;\theta)]^\top \hat\Omega_n [\mathbb{P}_n U(O;\theta)]$
for some positive-definite weighting matrix $\hat\Omega_n
\xrightarrow{p} \Omega$.

\begin{thm}{Asymptotic Distribution of the GMM Estimator}
\label{w12:thm:asym}
Suppose the IV moment condition holds, $O_i = (Y_i, T_i, X_i, Z_i)$
are i.i.d.\ with finite fourth moments, $A$ has full column rank, and
$\Sigma$ is positive definite.  Then
\[
  \sqrt{n}\,\bigl(\hat\theta_{\mathrm{GMM}} - \theta_0\bigr)
  \;\xrightarrow{d}\;
  N\bigl(0,\; V_{\mathrm{GMM}}(\Omega)\bigr),
\]
where
\begin{equation}
  V_{\mathrm{GMM}}(\Omega)
  = (A^\top\Omega A)^{-1}\,
    A^\top\Omega\,\Sigma\,\Omega A\,
    (A^\top\Omega A)^{-1}.
  \label{w12:eq:gmm-var}
\end{equation}
\end{thm}

\begin{tcolorbox}[remarkstyle, title={Intuition: the Sandwich Structure}]
The asymptotic variance has the same logic as every sandwich formula
in semiparametric estimation: one part measures how responsive the
moments are to the parameter ($A$), and another measures how variable
the moments themselves are ($\Sigma$).  Estimation is more precise
when the moments react strongly to the parameter (large $A$) and
fluctuate little across samples (small $\Sigma$).  The weighting
matrix $\Omega$ scales the two layers but does not alter their
fundamental roles.
\end{tcolorbox}

The formula~\eqref{w12:eq:gmm-var} is the \emph{sandwich variance}
for GMM, and is an instance of the general estimating-equation
theory from Chapter~9: the sensitivity matrix $A$ plays the role of
$-\E[\partial U/\partial\theta^\top]$ and the moment variance $\Sigma$
plays the role of $\E[UU^\top]$.

\subsection{Two Important Special Cases}

\textbf{Exactly identified case ($q = 1$, $m = k$).}
When there is a single instrument, $A$ is square and invertible by
the rank condition.  The GMM objective has an exact solution
$\mathbb{P}_n U(O;\hat\theta) = 0$ regardless of $\Omega$, so all
weighting matrices yield the same estimator.  The sandwich variance
simplifies to
\begin{equation}
  V_{\mathrm{IV}} = A^{-1}\,\Sigma\,(A^\top)^{-1},
  \label{w12:eq:iv-var}
\end{equation}
matching the general moment-estimator formula from Chapter~9.  In
the scalar no-covariate case ($p = 0$), $A = \E[ZT]$ and
$\Sigma = \E[\varepsilon^2 Z^2]$, so
\[
  V_{\mathrm{IV}}
  = \frac{\E[\varepsilon^2 Z^2]}{\bigl(\E[ZT]\bigr)^2}.
\]
Under homoskedasticity $\E[\varepsilon^2 \mid Z] = \sigma^2$, this
further reduces to $\sigma^2\,\mathrm{Var}(Z)^{-1}/\pi^2 =
\sigma^2/(\pi^2\,\mathrm{Var}(Z))$.

\textbf{Efficient GMM.}
Among all choices of $\Omega$, the asymptotic variance
$V_{\mathrm{GMM}}(\Omega)$ is minimized (in the positive
semidefinite sense) by $\Omega^\ast = \Sigma^{-1}$.  Substituting
into~\eqref{w12:eq:gmm-var}:
\begin{equation}
  V_{\mathrm{eff}}
  = \bigl(A^\top\Sigma^{-1}A\bigr)^{-1}.
  \label{w12:eq:eff-var}
\end{equation}
This is the \emph{efficient GMM variance} and is the semiparametric
efficiency lower bound for IV estimation within the linear
IV moment-restriction model defined by $\E[W\varepsilon] = 0$;
it is the best achievable variance among all estimators that use only
these moment conditions, but it does not in general coincide with
efficiency bounds from richer semiparametric models that exploit
additional structure.  For any other weighting matrix $\Omega$,
$V_{\mathrm{GMM}}(\Omega) - V_{\mathrm{eff}} \succeq 0$.

\begin{remark}[When 2SLS is efficient]
The 2SLS estimator corresponds to the weighting matrix
$\Omega_{\mathrm{2SLS}} = \E[W_i W_i^\top]^{-1}$.  Under
homoskedasticity, $\Sigma = \sigma^2\E[WW^\top]$, so
$\Sigma^{-1} = \sigma^{-2}\E[WW^\top]^{-1} \propto \Omega_{\mathrm{2SLS}}$.
Since scalar multiples of $\Omega$ do not change the minimizer,
2SLS achieves $V_{\mathrm{eff}}$ under homoskedasticity.  Under
heteroskedasticity, $\Sigma \neq \sigma^2\E[WW^\top]$ and 2SLS is
strictly less efficient than the optimal weighting
$\hat\Omega_n = (\mathbb{P}_n \hat\varepsilon_i^2 W_i W_i^\top)^{-1}$.
Efficient GMM improves over 2SLS only when the conditional variance
structure makes non-identity weighting informative; under
homoskedasticity, the optimal weighting collapses to the 2SLS case
and no efficiency gain is available.  This tells students when they
can stop at 2SLS without sacrificing asymptotic efficiency: whenever
homoskedasticity is plausible, efficient GMM brings no gain over 2SLS.
\end{remark}

\subsection{Consistent Variance Estimation}

Let $\hat\varepsilon_i = Y_i - D_i^\top\hat\theta$ denote the
structural residuals evaluated at the GMM estimates.  Consistent
estimators of $A$ and $\Sigma$ are
\[
  \hat{A} = \frac{1}{n}\sum_{i=1}^n W_i D_i^\top,
  \qquad
  \hat\Sigma = \frac{1}{n}\sum_{i=1}^n \hat\varepsilon_i^2\, W_i W_i^\top.
\]
Substituting into~\eqref{w12:eq:gmm-var} gives the heteroskedasticity-robust
sandwich variance estimator:
\begin{equation}
  \hat{V}_{\mathrm{GMM}}
  = (\hat{A}^\top\hat\Omega_n\hat{A})^{-1}\,
    \hat{A}^\top\hat\Omega_n\,\hat\Sigma\,\hat\Omega_n\hat{A}\,
    (\hat{A}^\top\hat\Omega_n\hat{A})^{-1}.
  \label{w12:eq:sandwich-est}
\end{equation}
In the exactly identified case, $\hat{A}$ is square and
\eqref{w12:eq:sandwich-est} reduces to
$\hat{A}^{-1}\hat\Sigma(\hat{A}^\top)^{-1}$, independent of
$\hat\Omega_n$.  For the efficient GMM estimator, the optimal
weighting is $\hat\Omega_n = \hat\Sigma^{-1}$, and the variance
estimator simplifies to $(\hat{A}^\top\hat\Sigma^{-1}\hat{A})^{-1}$.

In applications, the default variance estimator should rarely be
the i.i.d.\ formula.  Heteroskedasticity-robust, and when appropriate
cluster-robust, standard errors are the practical baseline for IV
work.  The cluster-robust version replaces $\hat\varepsilon_i^2
W_i W_i^\top$ in $\hat\Sigma$ with $(\sum_{i \in c}\hat\varepsilon_i
W_i)(\sum_{i \in c}\hat\varepsilon_i W_i)^\top$ summed over clusters
$c$, and is required whenever observations within a group share
unmodeled common shocks---as in the KIPP application of Chapter~7,
where test scores for students within the same lottery cohort may be
correlated.

\begin{remark}[Two-step efficient GMM]
Because the optimal weighting $\hat\Omega_n = \hat\Sigma^{-1}$
requires preliminary residuals $\hat\varepsilon_i$, efficient GMM
is typically implemented in two steps: obtain a consistent first-step
estimator (e.g., 2SLS) to form $\hat\varepsilon_i$, construct
$\hat\Sigma$, and then re-minimize the GMM criterion with
$\hat\Omega_n = \hat\Sigma^{-1}$.  The resulting estimator achieves
$V_{\mathrm{eff}}$ asymptotically.
\end{remark}


% ============================================================
%  Advanced enrichment begins here
% ============================================================
\begin{tcolorbox}[warnstyle, title={Advanced Enrichment:
  Sections~\ref{w12:sec:gel}--\ref{w12:sec:cf}}]
The remaining sections develop two extensions beyond the core
material.  They presuppose familiarity with
\S\S\,\ref{w12:sec:wald}--\ref{w12:sec:weak} but can be omitted
without loss of continuity.  Section~\ref{w12:sec:gel} introduces
generalized empirical likelihood (GEL) as a one-step alternative to
efficient GMM with better finite-sample properties.
Section~\ref{w12:sec:cf} develops the control function approach and
its extension to nonlinear triangular models.
\end{tcolorbox}

% ============================================================
\section{Generalized Empirical Likelihood}
\label{w12:sec:gel}
% ============================================================

\begin{remark}[Advanced topic]
This section introduces generalized empirical likelihood (GEL), an
alternative to GMM for estimating IV models from moment restrictions.
The main takeaway for this course is that every GEL estimator shares
the same first-order asymptotic efficiency as efficient GMM
(Theorem~\ref{w12:thm:gel-asymp}), while GEL can have better
higher-order finite-sample properties---in particular, smaller bias
in overidentified models with many instruments.  The detailed
duality formulation, the two-loop profile algorithm, and the
special-cases table are supplementary material and may be treated
as optional reading.
\end{remark}

Generalized empirical likelihood is a one-step alternative to
two-step efficient GMM.  It uses the same IV moment restrictions
but incorporates them through an implied reweighting of the sample
rather than through a separately estimated covariance weight matrix.
Its main attraction is that it attains the same first-order
efficiency as efficient GMM while sometimes improving finite-sample
behavior and avoiding the preliminary weighting step.

Section~\ref{w12:sec:asymp} showed that efficient GMM achieves the
semiparametric efficiency bound $(A^\top\Sigma^{-1}A)^{-1}$ but
requires a two-step procedure: a preliminary consistent estimator
is used to form $\hat\Sigma$, which is then plugged in as the
weighting matrix for the second-step optimization.  This two-step
structure introduces a practical question of how many iterations to
run and contributes a source of finite-sample bias
\citep{neweysmith2004}.  \emph{Generalized empirical likelihood} (GEL)
provides a one-step alternative: it jointly optimizes over the
structural parameter $\theta$ and an auxiliary dual variable,
implicitly constructing an optimal weighting without a separate
preliminary step.  All GEL estimators share the same first-order
asymptotic distribution as efficient GMM, making them
asymptotically equivalent alternatives that can offer better
finite-sample behavior.

\subsection{The GEL Estimator}
\label{w12:sec:gel:def}

Retain the notation of Section~\ref{w12:sec:asymp}: let
$U_i(\theta) = U(O_i;\theta) = W_i(Y_i - D_i^\top\theta)$ denote
the $m$-dimensional moment function, and write
$\bar{U}(\theta) = n^{-1}\sum_{i=1}^n U_i(\theta)$.

GEL introduces a strictly convex function
$G\colon \mathcal{V} \to \mathbb{R}$, where $\mathcal{V}$ is an
open interval containing zero.  The GEL estimator solves the
saddle-point problem
\begin{equation}
  \hat\theta_{\mathrm{GEL}}
  = \arg\min_{\theta}\;
    \sup_{\lambda \in \Lambda_n(\theta)}\;
    \frac{1}{n}\sum_{i=1}^n \bigl[-G\!\left(\lambda^\top U_i(\theta)\right)\bigr],
  \label{w12:eq:gel}
\end{equation}
where $\lambda \in \mathbb{R}^m$ is an auxiliary dual variable and
$\Lambda_n(\theta) = \{\lambda : \lambda^\top U_i(\theta) \in
\mathcal{V},\;i = 1,\ldots,n\}$ is the feasible set for $\lambda$
given $\theta$.  No preliminary weighting matrix is required: the
inner supremum over $\lambda$ performs the role of weighting
implicitly.

We normalize $G$ to satisfy three conditions:
\begin{equation}
  G(0) = 0,
  \qquad
  g(0) = 1,
  \qquad
  g'(0) = 1.
  \label{w12:eq:gel-norm}
\end{equation}
The conditions $g(0) = g'(0) = 1$ can always be imposed by a
linear rescaling of $G$ and $\lambda$ without changing the
estimator of $\theta$, and they ensure that the GEL first-order
conditions reduce to the efficient GMM conditions at the population
level.  The additional condition $G(0) = 0$ is a normalization of
the additive constant of $G$; it likewise has no effect on the
estimator, but it plays an essential role in the duality developed in
Section~\ref{w12:sec:gel:dual} by pinning down the correspondence
between $G$ and the primal divergence function.

The three most important special cases correspond to different
choices of $G$.

\begin{defn}{Special Cases of GEL}
\label{w12:defn:gel-cases}
\begin{enumerate}
  \item \textbf{Empirical likelihood (EL):}
    $G(v) = -\log(1 - v)$, $\mathcal{V} = (-\infty, 1)$.
  \item \textbf{Exponential tilting (ET):}
    $G(v) = e^v - 1$, $\mathcal{V} = \mathbb{R}$.
  \item \textbf{Continuous updating estimator (CUE):}
    $G(v) = v + v^2/2$, $\mathcal{V} = \mathbb{R}$.
    The function $G$ is quadratic, and for this choice the
    GEL saddle-point problem reduces to simultaneous minimization
    of the GMM criterion with the weighting matrix
    $\hat\Omega(\theta) = [n^{-1}\sum_i U_i(\theta) U_i(\theta)^\top]^{-1}$
    evaluated continuously at the current $\theta$.
\end{enumerate}
Each satisfies the normalization~\eqref{w12:eq:gel-norm}.
\end{defn}

\begin{thm}{CUE as a GEL Estimator \citep{neweysmith2004}}
\label{w12:thm:cue-gel}
If $G(v)$ is quadratic (i.e.\ $G(v) = v + v^2/2$ up to rescaling), then $\hat\theta_{\mathrm{GEL}} =
\hat\theta_{\mathrm{CUE}}$, where the continuous updating estimator
solves
\[
  \hat\theta_{\mathrm{CUE}}
  = \arg\min_\theta\;
    \bar{U}(\theta)^\top
    \Bigl[n^{-1}{\textstyle\sum_i} U_i(\theta)U_i(\theta)^\top\Bigr]^{-1}
    \bar{U}(\theta).
\]
\end{thm}

Unlike the two-step efficient GMM estimator, CUE uses a weighting
matrix that is continuously updated as $\theta$ changes, so no
preliminary estimator is needed.  It has the same first-order
asymptotics as efficient GMM but in practice requires a global
nonlinear optimizer because the criterion is no longer quadratic in
$\theta$.

\begin{tcolorbox}[remarkstyle, title={Optional Enrichment:
  Convex-Conjugate Duality (\S\ref{w12:sec:gel:dual})}]
The subsection below derives the GEL saddle-point problem as the
Lagrangian dual of a minimum-discrepancy primal problem.  This
duality is mathematically elegant and connects GEL to information
geometry, but it is not required for using GEL estimators in
practice.  Students who have covered the core GEL material above
may skip this subsection on a first pass.
\end{tcolorbox}

\subsection{The Convex-Conjugate Duality with Minimum Discrepancy}
\label{w12:sec:gel:dual}

The GEL saddle-point problem~\eqref{w12:eq:gel} is the Lagrangian
dual of a natural primal problem that re-weights observations rather
than moments.  The bridge between the two is the
Legendre--Fenchel conjugate of $G$.

\subsubsection*{The primal: minimum-discrepancy estimation}

The \emph{minimum-discrepancy} (MD) estimator minimizes a convex
divergence between a vector of observation weights
$\boldsymbol\omega = (\omega_1, \ldots, \omega_n)$ and the
reference weight $\omega_i = 1$ (one unit of mass per observation),
subject to the sample moment condition being exactly satisfied:
\begin{equation}
  \hat\theta_{\mathrm{MD}}
  = \arg\min_{\theta,\,\omega_1,\ldots,\omega_n}\;
    \sum_{i=1}^n F(\omega_i)
  \quad
  \text{subject to}
  \quad
  \sum_{i=1}^n \omega_i U_i(\theta) = 0,
  \label{w12:eq:md}
\end{equation}
where $F\colon(0,\infty)\to\mathbb{R}$ is a strictly convex function
with minimum at $\omega = 1$.  The problem asks: what is the
re-weighting of observations, closest to the uniform reference
$\omega_i = 1$ in the $F$-divergence sense, that exactly satisfies
the moment restriction?

\subsubsection*{The conjugate relationship}

The key structural fact is that the MD divergence $F$ and the
GEL strictly convex function $G$ are related through the Legendre--Fenchel
conjugate.  Define
\begin{equation}
  F(\omega)
  = \sup_{v \in \mathcal{V}}\,\bigl[\,\omega v - G(v)\bigr],
  \label{w12:eq:gel-conjugate}
\end{equation}
the Legendre--Fenchel conjugate of $G$.  Because $G$ is strictly
convex, the conjugate $F$ is itself strictly convex.  The normalization~\eqref{w12:eq:gel-norm} determines the
behavior of $F$ at the reference point $\omega = 1$: the supremum
in~\eqref{w12:eq:gel-conjugate} is attained at $v^* = 0$
(since the first-order condition gives $\omega = g(v^*)$, so
$\omega = 1 \Rightarrow v^* = 0$), yielding
\begin{equation}
  F(1) = 1 \cdot 0 - G(0) = 0,
  \qquad
  f(1) = v^*(1) = 0,
  \label{w12:eq:gel-F-norm}
\end{equation}
where the second equality uses the envelope theorem.  The location of
the minimizer follows entirely from $g(0) = 1$: the first-order
condition $\omega = g(v^*)$ gives $v^* = 0$ at $\omega = 1$,
and hence $f(1) = 0$, regardless of the value of $G(0)$.  The
condition $G(0) = 0$ has no bearing on which $\omega$ minimizes
$F$; it normalizes the \emph{level} of $F$ so that $F(1) = 0$,
making the divergence interpretable as a non-negative cost with zero
cost at the reference weight $\omega_i = 1$.

\subsubsection*{The duality theorem}

\begin{thm}{GEL--MD Duality \citep{neweysmith2004,ragusa2011}}
\label{w12:thm:gel-md}
Let $G$ satisfy the normalization~\eqref{w12:eq:gel-norm} and
let $F$ be defined by~\eqref{w12:eq:gel-conjugate}.  Then the GEL
problem~\eqref{w12:eq:gel} is the Lagrangian dual of the MD
problem~\eqref{w12:eq:md}: their first-order conditions coincide, so
$\hat\theta_{\mathrm{GEL}} = \hat\theta_{\mathrm{MD}}$.  The dual
variable $\hat\lambda$ and the primal weights $\hat\omega_i$ are
related by
\[
  \hat\omega_i = g\!\bigl(\hat\lambda^\top U_i(\hat\theta)\bigr),
\]
which satisfies $f(\hat\omega_i) = \hat\lambda^\top U_i(\hat\theta)$
by the conjugate relationship $f^{-1} = g$.
\end{thm}

\citet{neweysmith2004} established this duality for the
\emph{Cressie--Read family} (Cressie and Read, 1984), in which
$G(v) = (1+\gamma v)^{(\gamma+1)/\gamma}/(\gamma+1) -
1/(\gamma+1)$ and $F(\omega) = [(\omega)^{\gamma+1} - 1 - (\gamma+1)
(\omega-1)]/[\gamma(\gamma+1)]$.  EL, ET, and CUE correspond to
$\gamma = -1$, $\gamma = 0$, and $\gamma = 1$.  \citet{ragusa2011}
showed that the duality holds in full generality for any $G$
satisfying~\eqref{w12:eq:gel-norm}: the Cressie--Read parametrization
is not needed for the structural result; the convex conjugate is the
correct unifying framework.

\subsubsection*{Solving the MD Problem: Two-Loop Profile Optimization}

Although the MD problem~\eqref{w12:eq:md} jointly optimizes over
$(\theta, \omega_1,\ldots,\omega_n)$, it can be solved by a
two-loop profile procedure that exploits the convexity of $F$ in
the weights.

\textbf{Inner loop: profile weights at fixed $\theta$.}  For a
given $\theta$, the minimization over $(\omega_1,\ldots,\omega_n)$
subject to $\sum_i \omega_i U_i(\theta) = 0$ is a convex program.
Its Lagrangian is
\[
  L(\omega, \lambda)
  = \sum_{i=1}^n F(\omega_i)
    - \lambda^\top \sum_{i=1}^n \omega_i U_i(\theta),
\]
and the first-order condition $f(\omega_i) = \lambda^\top
U_i(\theta)$ gives the profile weight
\[
  \hat\omega_i(\theta) = f^{-1}\!\bigl(\lambda^\top
  U_i(\theta)\bigr) = g\!\bigl(\lambda^\top U_i(\theta)\bigr),
\]
where the last equality uses the conjugate identity $f^{-1} = g$.
Substituting into the moment constraint $\sum_i \hat\omega_i(\theta)
U_i(\theta) = 0$ determines the Lagrange multiplier
$\hat\lambda(\theta)$ as the solution to
\begin{equation}
  \hat\lambda(\theta)
  = \arg\min_{\lambda \in \Lambda_n(\theta)}\;
    \frac{1}{n}\sum_{i=1}^n G\!\bigl(\lambda^\top U_i(\theta)\bigr).
  \label{w12:eq:gel-inner}
\end{equation}
This inner minimization is a convex optimization in $\lambda$ alone
and can be solved efficiently by Newton's method for each fixed
$\theta$.

\textbf{Profile divergence.}  Once $\hat\omega_i(\theta) =
g(\hat\lambda(\theta)^\top U_i(\theta))$ is substituted back into the
MD objective, the Fenchel--Young equality $F(\omega) = \omega
f(\omega) - G(f(\omega))$ and the moment constraint
$\sum_i \hat\omega_i(\theta) U_i(\theta) = 0$ together give
\begin{equation}
  Q_{\mathrm{MD}}(\theta)
  \;\equiv\;
  \sum_{i=1}^n F\!\bigl(\hat\omega_i(\theta)\bigr)
  = -\sum_{i=1}^n G\!\bigl(\hat\lambda(\theta)^\top U_i(\theta)\bigr).
  \label{w12:eq:md-profile}
\end{equation}
The profiled divergence therefore reduces to a function of $\theta$
alone, evaluated at the inner-loop solution.

\textbf{Outer loop: minimize over $\theta$.}  The MD estimator is
\[
  \hat\theta_{\mathrm{MD}}
  = \arg\min_{\theta}\; Q_{\mathrm{MD}}(\theta)
  = \arg\max_{\theta}\;
    \min_{\lambda \in \Lambda_n(\theta)}\;
    \frac{1}{n}\sum_{i=1}^n G\!\bigl(\lambda^\top U_i(\theta)\bigr),
\]
where the second equality follows from~\eqref{w12:eq:md-profile}.
This is exactly the GEL saddle-point problem~\eqref{w12:eq:gel}
written as a max--min, confirming at the algorithmic level that the
two procedures find the same estimator.  In practice the outer loop
iterates over $\theta$, calling the inner loop at each step to
evaluate $\hat\lambda(\theta)$ and the gradient of
$Q_{\mathrm{MD}}(\theta)$.

\begin{remark}[Connection to the GEL saddle-point]
The two-loop MD procedure and the GEL saddle-point are two
computational routes to the same estimator.  The GEL formulation
jointly optimizes $(\theta, \lambda)$ in the saddle-point
problem~\eqref{w12:eq:gel}, while the MD two-loop approach decouples
the problem: the inner loop solves a convex minimization in $\lambda$
for fixed $\theta$, and the outer loop minimizes the profiled
divergence in $\theta$.  The inner-loop solution $\hat\lambda(\theta)$
from~\eqref{w12:eq:gel-inner} is precisely the dual variable that
appears in the GEL first-order conditions, and the profile
divergence~\eqref{w12:eq:md-profile} is the GEL profile objective.
\end{remark}

\subsubsection*{Empirical probabilities}

The primal weights $\hat\omega_i$ need not sum to one.  Normalizing
gives the \emph{empirical probabilities}
\begin{equation}
  \hat\pi_i
  = \frac{\hat\omega_i}{\sum_{j=1}^n \hat\omega_j}
  = \frac{g\!\left(\hat\lambda^\top U_i(\hat\theta)\right)}
         {\displaystyle\sum_{j=1}^n
          g\!\left(\hat\lambda^\top U_j(\hat\theta)\right)},
  \qquad i = 1,\ldots,n,
  \label{w12:eq:gel-probs}
\end{equation}
where $\hat\lambda =
\arg\inf_{\lambda \in \Lambda_n(\hat\theta)}
n^{-1}\sum_i G(\lambda^\top U_i(\hat\theta))$.
These satisfy $\sum_i \hat\pi_i = 1$ and $\sum_i \hat\pi_i
U_i(\hat\theta) = 0$ at the GEL first-order conditions.  For CUE,
the primal weights are generally non-uniform and are induced by the
quadratic function $G$.

\subsubsection*{Special cases and divergence functions}

The conjugate relationship~\eqref{w12:eq:gel-conjugate} determines
the MD divergence $F$ from $G$.  Table~\ref{w12:tab:gel-cases}
summarizes the three main special cases; the derivations follow by
solving the first-order condition $\omega - g(v^*) = 0$ for $v^*$
and substituting back into~\eqref{w12:eq:gel-conjugate}.

\begin{table}[ht]
\centering
\small
\begin{tabular}{@{}llccc@{}}
\toprule
Estimator & $G(v)$ & $g(v) = G'(v)$ & $F(\omega)$ & Divergence \\
\midrule
EL  & $-\log(1-v)$ & $\displaystyle\frac{1}{1-v}$
    & $\omega - 1 - \log\omega$
    & KL: $\sum_i \hat\pi_i \log\frac{\hat\pi_i}{1/n}$ \\[8pt]
ET  & $e^v - 1$ & $e^{v}$
    & $\omega\log\omega - \omega + 1$
    & Reverse KL: $\sum_i \frac{1}{n}\log\frac{1/n}{\hat\pi_i}$ \\[6pt]
Hellinger & $\displaystyle\frac{2v}{2-v}$ & $\displaystyle\frac{4}{(2-v)^2}$
    & $2(\sqrt{\omega}-1)^2$
    & $2n\sum_i\!\bigl(\sqrt{\hat\pi_i}-\tfrac{1}{\sqrt{n}}\bigr)^{\!2}$ \\[8pt]
CUE & $v + \tfrac{1}{2}v^2$ & $1 + v$
    & $\tfrac{1}{2}(\omega - 1)^2$
    & $\chi^2$: $\tfrac{n}{2}\sum_i (\hat\pi_i - \tfrac{1}{n})^2$ \\
\bottomrule
\end{tabular}
\caption{The four main GEL special cases.  $G$ is the strictly
  convex function defining the GEL estimator~\eqref{w12:eq:gel};
  $g = G'$ appears in the empirical probability
  formula~\eqref{w12:eq:gel-probs} via $\hat\omega_i = g(\hat\lambda^\top
  U_i(\hat\theta))$; $F(\omega) = \sup_{v}\,[\omega v - G(v)]$ is
  the Legendre--Fenchel conjugate of $G$ and defines the MD
  primal~\eqref{w12:eq:md}.  All four satisfy $G(0)=0$,
  $g(0)=1$, $F(1)=0$, $f(1)=0$.  The divergence column shows
  $\sum_i F(\hat\omega_i)$ expressed in terms of the normalized
  probabilities $\hat\pi_i = \hat\omega_i/\sum_j\hat\omega_j$.
  EL, Hellinger, ET, and CUE correspond to $\gamma=-1$,
  $\gamma=-\tfrac{1}{2}$, $\gamma=0$, and $\gamma=1$ in the
  Cressie--Read family, respectively.  The domain of $G$ is
  $\mathcal{V}=(-\infty,1)$ for EL, $\mathcal{V}=(-\infty,2)$ for
  Hellinger, and $\mathcal{V}=\mathbb{R}$ for ET and CUE.}
\label{w12:tab:gel-cases}
\end{table}

For EL, $F(\omega) = \omega - 1 - \log\omega$ is the
Kullback--Leibler divergence from $\hat\pi$ to the uniform $1/n$,
so the EL weights are the nonparametric maximum likelihood estimate
of the observation probabilities subject to the moment constraint.
For ET, the roles of the two distributions in the KL divergence are
reversed.  For CUE, the divergence is the squared Euclidean
distance from the reference weight.

\begin{remark}[Re-weighting observations vs.\ re-weighting moments]
The convex-conjugate duality highlights a structural contrast between
GMM and GEL.  GMM re-weights the \emph{moment conditions} via the
matrix $\hat\Omega$, assigning more importance to moment conditions
that are precisely estimated.  MD---and its GEL dual---instead
re-weights the \emph{observations}, finding the empirical
distribution closest to uniform in the $F$-divergence sense that is
consistent with the moment restrictions.  The two approaches use the
overidentification to improve efficiency through structurally
different mechanisms, yet they are asymptotically equivalent to
first order.
\end{remark}

\subsection{Asymptotic Properties and Comparison with GMM}
\label{w12:sec:gel:asymp}

Under the same regularity conditions as Theorem~\ref{w12:thm:asym},
every GEL estimator achieves the efficient GMM variance.

\begin{thm}{First-Order Asymptotic Equivalence \citep{neweysmith2004}}
\label{w12:thm:gel-asymp}
Suppose the moment condition $\E[U(O;\theta_0)] = 0$ holds, $O_i$
are i.i.d.\ with sufficient moments, the rank condition holds, and
$G$ satisfies the regularity conditions of
\citet{neweysmith2004}.  Then for any GEL estimator,
\[
  \sqrt{n}\,\bigl(\hat\theta_{\mathrm{GEL}} - \theta_0\bigr)
  \;\xrightarrow{d}\;
  N\!\left(0,\; (A^\top\Sigma^{-1}A)^{-1}\right),
\]
the same limiting distribution as the efficient GMM estimator.
\end{thm}

The first-order equivalence means that the choice among GEL variants,
or between GEL and efficient GMM, does not matter in large samples.
Differences emerge at higher order.

\begin{thm}{GEL Overidentification Test (Wilks' Theorem for GEL)
  \citep{neweysmith2004}}
\label{w12:thm:gel-wilks}
Under the same conditions as Theorem~\ref{w12:thm:gel-asymp}, the
GEL profile divergence evaluated at the estimator satisfies
\begin{equation}
  T_{\mathrm{GEL}}
  \;\equiv\;
  2\sum_{i=1}^n F\!\bigl(\hat\omega_i\bigr)
  \;=\;
  -2\sum_{i=1}^n G\!\bigl(\hat\lambda^\top U_i(\hat\theta_{\mathrm{GEL}})\bigr)
  \;\xrightarrow{d}\;
  \chi^2_{m-k},
  \label{w12:eq:gel-wilks}
\end{equation}
where $m$ is the number of moment conditions, $k$ is the number of
parameters, and $m - k$ is the number of overidentifying
restrictions.
\end{thm}

The two expressions in~\eqref{w12:eq:gel-wilks} are equal by the
profile divergence identity~\eqref{w12:eq:md-profile}: the total
primal divergence from the reference weights equals minus the
profiled GEL objective at the solution.  The $\chi^2_{m-k}$ limit
parallels Wilks' theorem for parametric likelihood ratio tests: the
degrees of freedom count the overidentifying restrictions, and the
statistic is zero under exact identification ($m = k$).  This is the
GEL counterpart of the Sargan--Hansen $J$-statistic from
Section~\ref{w12:sec:gmm}.

For EL specifically, $F(\omega) = \omega - 1 - \log\omega$, so
\[
  T_{\mathrm{EL}}
  = 2\sum_{i=1}^n \hat\omega_i \log\!\left(\frac{\hat\omega_i}{1}\right)
  \;=\; 2n\sum_{i=1}^n \hat\pi_i \log(n\hat\pi_i),
\]
the empirical likelihood ratio statistic of Qin and Lawless (1994),
which coincides with $-2$ times the log empirical likelihood ratio
under the moment constraint relative to the unconstrained maximum.

\begin{remark}[Comparison with the Sargan--Hansen $J$-statistic]
The $J$-statistic of Section~\ref{w12:sec:gmm} equals
$n\,\bar{U}(\hat\theta)^\top \hat\Sigma^{-1} \bar{U}(\hat\theta)$
and also converges to $\chi^2_{m-k}$ under the null.  The two
statistics test the same null hypothesis---all $m$ moment conditions
hold simultaneously---but differ in construction.  The $J$-statistic
uses the sample mean $\bar{U}(\hat\theta)$ and a separate estimate
$\hat\Sigma$, while $T_{\mathrm{GEL}}$ uses the profile divergence
at the re-weighted distribution.  For the CUE, $F(\omega) =
(\omega-1)^2/2$, so $T_{\mathrm{CUE}} = \sum_i (\hat\omega_i - 1)^2$,
a quadratic form that is closely related to the $J$-statistic evaluated
at the continuously updated weighting matrix.  Under the null all GEL
statistics share the same $\chi^2_{m-k}$ limit; they differ in local
power and higher-order properties.
\end{remark}

\begin{remark}[Finite-sample advantages of GEL over two-step GMM]
\citet{neweysmith2004} show that, to higher asymptotic order, GEL
estimators have smaller bias than standard two-step GMM in two
important respects.  First, GEL eliminates the bias component arising
from estimation of the Jacobian
$\E[\partial U_i(\theta_0)/\partial\theta^\top]$.  This is an
important source of finite-sample bias in overidentified IV models,
and it appears in GMM because the gradient term in the first-order
conditions is estimated by a sample average rather than by an
efficient weighted average.  The inner saddle-point over $\lambda$ in
GEL implicitly uses an efficient estimator of the Jacobian, so this
bias component vanishes.  Second, EL additionally eliminates the bias
from estimating the weighting matrix $\Sigma$, because EL uses the
re-weighted probabilities $\hat\pi_i$ as an efficient estimate of the
second moment.  Under symmetric disturbances---a condition satisfied
in many IV settings---these bias reductions extend to the full GEL
family.  In practice, these properties make GEL estimators attractive
in overidentified IV settings with moderate sample sizes or many
instruments, where finite-sample bias is a genuine concern.
\end{remark}

\begin{remark}[Computational considerations]
The GEL saddle-point problem~\eqref{w12:eq:gel} is a joint nonlinear
optimization in $(\theta, \lambda)$.  For the linear IV model, the
inner maximization over $\lambda$ can often be solved efficiently by
iterative algorithms, and standard software provides reliable GEL
implementations.  Among the three named special cases, CUE is the
most transparent computationally---the inner supremum has a closed
form for fixed $\theta$---but its objective is non-convex in $\theta$
and requires a global optimizer.  EL is often preferred in practice
for its minimum-bias properties; it additionally supports empirical
likelihood confidence regions that have favorable coverage properties
without requiring symmetry assumptions on the moment conditions.
\end{remark}


% ============================================================
\section{Weak Instruments and Inferential Fragility}
\label{w12:sec:weak}
% ============================================================

The GMM and 2SLS estimators are well-behaved when the instrument
strongly predicts treatment.  When this first-stage relationship is
weak, IV estimation becomes severely fragile.  A weak instrument is
not merely an efficiency problem: it makes the finite-sample
distribution of IV estimators highly non-normal, magnifies bias
toward OLS, and undermines conventional confidence intervals.  These
are not asymptotic curiosities but practical hazards that arise
routinely in empirical work, and they are among the most important
practical concerns in IV design.

\subsection{The Weak-Instrument Problem}

The 2SLS closed form~\eqref{w12:eq:2sls-closed} divides by the
sample covariance $\widehat{\mathrm{Cov}}(\tilde{Z}, T)$.  When the
population first-stage coefficient $\pi$ is near zero, this denominator
is near zero in the population and highly variable in finite samples.
The consequences are well-documented
\citep{bound1995problems}:
\begin{itemize}
  \item \textbf{Finite-sample bias toward OLS.}  The 2SLS estimator is
    biased in finite samples even under valid identification.  As
    $\pi \to 0$, the bias approaches that of OLS rather than zero.
  \item \textbf{Non-Gaussian finite-sample distribution.}  The
    distribution of $\hat\beta_{\mathrm{2SLS}}$ can be highly skewed
    or heavy-tailed, even in large samples, rendering the
    $N(0,\,V_{\mathrm{2SLS}})$ approximation unreliable.
  \item \textbf{Size distortion.}  Wald-type confidence intervals
    derived from the asymptotic normal approximation can severely
    undercover the true parameter.
\end{itemize}

\subsection{Diagnostic: The First-Stage $F$-Statistic}

The most widely used diagnostic for weak instruments is the
$F$-statistic from the first-stage regression, testing the joint
significance of $Z$ after partialling out $X$.  Staiger and Stock
(1997) argued informally for the rule of thumb $F \geq 10$ as an
indication of adequate instrument strength.  Stock and Yogo (2005)
provided formal critical values, tabulating thresholds at which
the 2SLS bias relative to OLS bias exceeds a specified tolerance.
The textbook first-stage $F$-statistic is most directly
interpretable in the single-endogenous-regressor setting; with
heteroskedasticity, clustering, multiple endogenous regressors, or
many instruments, more specialized diagnostics are preferable.
In all cases, a large first-stage $F$ supports relevance, but says
nothing directly about exogeneity or exclusion: a strong but invalid
instrument passes the first-stage test with no difficulty.  The
first-stage $F$-statistic is a relevance diagnostic, not a
certificate of instrument validity.

\begin{warn}{The $F > 10$ Rule of Thumb}
The threshold $F_{\text{first stage}} > 10$ is a widely cited but
coarse diagnostic.  It targets finite-sample bias relative to OLS,
not validity of the exclusion restriction.  It may also be too lenient
when there are multiple endogenous regressors or many instruments.
A more reliable diagnostic is the effective $F$-statistic of Olea and
Pflueger (2013), which remains valid under heteroskedasticity and
within-cluster correlation.
\end{warn}

\subsection{Weak-Instrument-Robust Inference}

When first-stage strength is uncertain, the standard normal
approximation for $\hat\beta_{\mathrm{2SLS}}$ may be unreliable, and
alternative inferential procedures with size guarantees are needed.

The \textbf{Anderson--Rubin (AR) test} (Anderson and Rubin, 1949)
inverts the question: rather than asking whether the data are
consistent with a specific estimate, it tests whether a hypothesized
value $\beta_0$ is consistent with the IV moment condition.
Substituting $\beta_0$ into the structural model gives
$Y - \beta_0 T = \alpha + \gamma^\top X + (\varepsilon + (\beta - \beta_0) T)$.
If $\beta_0 = \beta$, the composite error $\varepsilon + (\beta - \beta_0) T$
is uncorrelated with $Z$ by instrument validity.
The AR test regresses $Y - \beta_0 T$ on $Z$ and $X$, and tests the
null that the coefficient on $Z$ is zero.  Under homoskedasticity,
the resulting $F$-statistic has an exact $F$-distribution regardless
of instrument strength.  Inverting this test over a grid of $\beta_0$
values yields a confidence set that is valid whether or not the
instrument is weak.

The \textbf{conditional likelihood ratio (CLR) test} of Moreira
(2003) extends the AR idea more efficiently to the multiple-instrument
case.  In the single-endogenous-regressor, single-instrument benchmark,
the practical distinction between AR and CLR is much smaller; with
multiple instruments, CLR can be substantially more powerful.

\begin{remark}[Two distinct IV concerns]
The weak-instrument problem is a \emph{statistical} concern: it
questions the quality of estimation and inference given a valid
identification strategy.  It is distinct from the \emph{interpretive}
concern raised in Chapter~7---that even a perfectly strong instrument
identifies only the LATE for compliers, not the population ATE.
Any complete IV analysis must address both.
\end{remark}


% ============================================================
\section{The Control Function Approach}
\label{w12:sec:cf}
% ============================================================

The control-function approach offers another route to handling
endogeneity: instead of projecting treatment onto instruments and
imposing orthogonality in moment form, it augments the outcome model
with a control variable that absorbs the endogenous component of
treatment selection.  The endogenous part of treatment is captured
by the first-stage residual; including it in the outcome regression
controls away the correlation between treatment and the structural
error, restoring consistency.

Two-stage least squares achieves identification by replacing the endogenous
regressor $T$ with its exogenous projection $\hat{T}$.  The \emph{control
function} approach achieves the same goal by a different route: it adds the
first-stage residual to the outcome regression as an explicit control for
the endogenous variation, rather than removing it from the treatment variable.
In the linear model the two approaches are numerically equivalent;
Section~\ref{w12:sec:cf:linear} develops this equivalence and the associated
endogeneity test.  Section~\ref{w12:sec:cf:nonlinear} closes with a brief
note on what changes in nonlinear outcome models.

\subsection{Linear Model and Equivalence to 2SLS}
\label{w12:sec:cf:linear}

Return to the structural model
\begin{align*}
  Y &= \alpha + \beta T + \gamma^\top X + \varepsilon, \\
  T &= \pi^\top Z + \delta^\top X + \eta,
\end{align*}
where $\mathrm{Cov}(\varepsilon, \eta) \neq 0$ is the source of
endogeneity.  Write the structural error as
\[
  \varepsilon = \rho\,\eta + \xi,
  \qquad
  \E[\xi \mid \eta, Z, X] = 0,
\]
where $\rho = \mathrm{Cov}(\varepsilon, \eta)/\mathrm{Var}(\eta)$ is the
linear projection coefficient and $\xi$ is the orthogonal component of
$\varepsilon$ with respect to $\eta$.  Substituting into the outcome
equation:
\[
  Y = \alpha + \beta T + \gamma^\top X + \rho\,\eta + \xi.
\]
Because $\E[\xi \mid T, X, \eta] = 0$, conditioning on $\eta$ (or any
sufficient statistic for it) renders $T$ exogenous in the augmented
regression.  This motivates the following two-step procedure.

\begin{tcolorbox}[defstyle, title={Control Function Estimator (Linear Case)}]
\label{w12:box:cf-linear}
\begin{enumerate}
  \item \textbf{First stage.}  Regress $T$ on $Z$ and $X$ by OLS; obtain
    residuals $\hat\eta_i = T_i - \hat\pi^\top Z_i - \hat\delta^\top X_i$.
  \item \textbf{Second stage.}  Regress $Y$ on $T$, $X$, and $\hat\eta$
    by OLS; the coefficient on $T$ is the control function estimator
    $\hat\beta_{\mathrm{CF}}$.
\end{enumerate}
\end{tcolorbox}

\begin{thm}{Equivalence of 2SLS and the Control Function Estimator}
\label{w12:thm:cf-equiv}
In the linear IV model with instruments $Z$ and covariates $X$, the
control function estimator $\hat\beta_{\mathrm{CF}}$ from the augmented
OLS regression of $Y$ on $(T, X, \hat\eta)$ is numerically identical to
the 2SLS estimator $\hat\beta_{\mathrm{2SLS}}$.
\end{thm}

\begin{proof}
The Frisch--Waugh--Lovell theorem states that the coefficient on $T$ in
the OLS regression of $Y$ on $(T, X, \hat\eta)$ equals the coefficient
on $\tilde{T}$ in the regression of $\tilde{Y}$ on $\tilde{T}$, where
tildes denote residuals after projecting out $(X, \hat\eta)$.  Since
$\hat\eta = T - \hat{T}$, projecting out $\hat\eta$ and $X$ from $T$
is equivalent to projecting out $X$ and $T - \hat{T}$ from $T$, which
leaves $\hat{T}$ after partialling out $X$.  Thus $\tilde{T}$ is the
residual from projecting $\hat{T}$ on $X$, which is the same residual
used in the 2SLS second stage.  The coefficient is therefore identical
to $\hat\beta_{\mathrm{2SLS}}$.
\end{proof}

\begin{remark}[Standard errors in the second stage]
Although the coefficient on $T$ is identical to 2SLS, the OLS standard
errors from the augmented regression are not.  They ignore the sampling
variation in the estimated residuals $\hat\eta_i$ and are therefore
incorrect.  Correct inference requires either the 2SLS sandwich formula
from Section~\ref{w12:sec:asymp}, or a bootstrap that resamples both
stages jointly.
\end{remark}

\subsection{Testing Endogeneity via Residual Inclusion}

The control function representation yields a natural test of the null
hypothesis $H_0\colon \rho = 0$, which is the hypothesis that $T$ is
exogenous.  Under $H_0$, the structural error $\varepsilon$ is
uncorrelated with $\eta$, so $\hat\eta$ should have no predictive power
for $Y$ after controlling for $T$ and $X$.  The $t$-statistic on
$\hat\eta$ in the augmented regression is therefore a test of
endogeneity.  A rejection suggests endogeneity under the maintained
instrument assumptions; it is not a stand-alone validation of the
instrument itself, since the test takes instrument validity as given.

\begin{remark}[Hausman test connection]
This $t$-test on the first-stage residual is the regression-based
form of the Hausman (1978) endogeneity test.  The test compares the
OLS and 2SLS estimates of $\beta$: a large discrepancy indicates
that the OLS bias from $\mathrm{Cov}(T, \varepsilon) \neq 0$ is
substantial.  The control function formulation makes the mechanics
of this comparison explicit.
\end{remark}

\subsection{Brief Note on Nonlinear Extensions}
\label{w12:sec:cf:nonlinear}

The equivalence of §\ref{w12:sec:cf:linear} is a theorem about the
\emph{linear additive} IV model.  Once the outcome equation is nonlinear,
the situation changes substantially.

\textbf{2SLS does not carry over.}  In a probit or Poisson outcome
model, for instance, 2SLS is generally inconsistent: plugging in
$\hat T$ breaks the nonlinear link function, and the moment condition
$\E[Z\varepsilon] = 0$ does not translate into the likelihood.

\textbf{Related control-variable methods exist.}  For certain
nonlinear triangular models, a control-variable strategy can recover
conditional exogeneity, but the control variable is not the first-stage
OLS residual $\hat\eta$.  \citet{imbens2009} showed that under two
additional conditions---\emph{independence}, $(\varepsilon, \eta) \indep Z$,
and \emph{scalar monotonicity}, $T = h(Z, \eta)$ strictly monotonic in a
scalar $\eta$---the conditional CDF
\begin{equation}
  V \;=\; F_{T \mid Z, X}(T \mid Z, X)
  \label{w12:eq:cf-control-var}
\end{equation}
is a valid control variable: $T \indep \varepsilon \mid V$.
Conditioning on $V$ recovers structural variation in $T$, allowing
identification of the average structural function and related estimands
by a two-step nonparametric procedure.

\textbf{Substantially more structure is required.}  The independence
and scalar-monotonicity conditions are much stronger than the linear-model
IV assumptions, and they rule out many simultaneous-equation settings
with multiple disturbances.  Estimation requires nonparametric
first- and second-stage regressions, and convergence rates are slower
than in the linear case.  These methods are beyond the scope of this
chapter; \citet{imbens2009} is the primary reference.

\begin{remark}[Control function vs.\ 2SLS: when they differ]
In the linear additive model, the control function estimator and 2SLS
are identical (Theorem~\ref{w12:thm:cf-equiv}).  They diverge as soon as
the outcome equation is nonlinear: 2SLS is generally inconsistent there,
while the Imbens--Newey control-variable approach may remain valid under
the stronger assumptions above.  The added flexibility of control-function
methods comes at the price of stronger modeling structure, especially in
nonlinear settings.
\end{remark}



% ============================================================
\section{Chapter Summary}
\label{w12:sec:summary}
% ============================================================
This chapter developed the main estimation tools for the linear
instrumental-variable model.  The organizing principle is the same as
in Chapter~9: identification provides a target, and the estimator is
a sample analog of the identification formula.  The key results are:

\begin{enumerate}
  \item \textbf{Wald estimator, reduced form, and first stage.}
    The \textbf{Wald estimator} is the direct sample analog of the
    Wald estimand for a binary instrument with no covariates: the
    ratio of the reduced-form difference-in-means to the first-stage
    difference-in-means.  The first stage answers how strongly the
    instrument moves treatment; the reduced form answers how the
    instrument changes the outcome; the ratio converts the
    intent-to-treat effect into a per-unit-of-treatment causal effect
    under the IV assumptions.

  \item \textbf{IV regression estimator and 2SLS.}
    The \textbf{IV regression estimator} extends the Wald estimator
    to the general single-instrument setting with covariates $X$, as
    the ratio of within-$X$ instrument-outcome to
    instrument-treatment sample covariances, derived via the
    constrained normal equations.  \textbf{Two-stage least squares}
    extends this further to multiple instruments: the first stage
    projects $T$ onto $(Z, X)$ and the second stage regresses $Y$
    on the fitted values and $X$.

  \item \textbf{Equivalence and moment conditions.}
    The IV regression estimator and 2SLS are \textbf{numerically
    identical} in the single-instrument case for any $X$
    (Theorem~\ref{w12:thm:equiv}), with the Wald estimator as the
    further special case $Z \in \{0,1\}$, $X$ absent.  Both are
    \textbf{method-of-moments estimators} based on the IV
    orthogonality condition $\E[W\varepsilon] = 0$, an instance of
    the general estimating-equation framework of Chapter~9.

  \item \textbf{GMM and efficient weighting.}
    In overidentified models, \textbf{efficient GMM} generalizes
    2SLS and achieves the semiparametric efficiency bound within the
    maintained IV moment-restriction model.  2SLS achieves that bound
    only under homoskedasticity; under heteroskedasticity, efficient
    GMM strictly dominates.

  \item \textbf{Robust variance estimation.}
    The \textbf{asymptotic distribution} of the GMM estimator is
    normal with sandwich variance
    $V_{\mathrm{GMM}} = (A^\top\Omega A)^{-1} A^\top\Omega\Sigma\Omega A
    (A^\top\Omega A)^{-1}$, where $A$ is the sensitivity matrix and
    $\Sigma$ is the moment variance.  In applications, the default
    should be heteroskedasticity-robust standard errors; cluster-robust
    standard errors are required when observations within groups share
    common shocks.

  \item \textbf{Weak instruments and inferential fragility.}
    \textbf{Weak instruments} cause finite-sample bias toward OLS,
    non-Gaussian distributions, and size distortion in conventional
    Wald confidence intervals.  The first-stage $F$-statistic is a
    relevance diagnostic, not a certificate of instrument validity.
    Anderson--Rubin confidence sets provide inference robust to
    weak instruments regardless of instrument strength.

  \item \textbf{GEL as an advanced extension.}
    \textbf{Generalized empirical likelihood} is a family of one-step
    estimators that achieve the efficient GMM variance without a
    preliminary weighting step (Theorem~\ref{w12:thm:gel-asymp}).
    Key special cases are empirical likelihood (EL), exponential
    tilting (ET), and the continuous updating estimator (CUE).  The
    GEL saddle-point problem is the Lagrangian dual of a
    minimum-discrepancy primal problem that re-weights observations
    rather than moments.  To higher asymptotic order, GEL eliminates
    the Jacobian-estimation bias of two-step GMM.

  \item \textbf{Control function as an alternative modeling route.}
    In the linear model, the \textbf{control-function approach} is
    numerically equivalent to 2SLS
    (Theorem~\ref{w12:thm:cf-equiv}); it augments the outcome
    regression with the first-stage residual rather than replacing
    the endogenous regressor.  Beyond the linear model, related
    control-variable methods exist for nonlinear triangular models
    under the stronger independence and scalar-monotonicity
    conditions of \citet{imbens2009}.  The added flexibility of
    control-function methods in nonlinear settings comes at the price
    of stronger modeling structure.
\end{enumerate}

% ============================================================
\section*{Exercises}
\label{w12:sec:exercises}
% ============================================================

\begin{enumerate}

% ------------------------------------------------------------------
\item \textbf{The Wald estimator as a ratio-of-moments estimator.}
  Consider the binary-instrument, no-covariate setting of
  Section~\ref{w12:sec:wald}.  Let $\mu_Y(z) = \E[Y \mid Z = z]$
  and $\mu_T(z) = \E[T \mid Z = z]$ for $z \in \{0,1\}$, and write
  $\Delta_Y = \mu_Y(1) - \mu_Y(0)$, $\Delta_T = \mu_T(1) - \mu_T(0)$.
  \begin{enumerate}
    \item Express the Wald estimand $\beta = \Delta_Y / \Delta_T$ as
      the solution to a pair of moment conditions of the form
      $\E[U(O;\beta)] = 0$, where $U \in \mathbb{R}^2$.  Identify $U$
      explicitly and verify that the system is exactly identified.
    \item Using the fact that $\bar{Y}_z \xrightarrow{p} \mu_Y(z)$ and
      $\bar{T}_z \xrightarrow{p} \mu_T(z)$, prove that
      $\hat\beta_{\mathrm{Wald}} \xrightarrow{p} \beta$ via the
      continuous mapping theorem.  State the condition on $\Delta_T$
      that is required.
    \item Apply the delta method to the vector
      $(\hat\Delta_Y, \hat\Delta_T)^\top$ to show that
      \[
        \sqrt{n}\,(\hat\beta_{\mathrm{Wald}} - \beta)
        \;\xrightarrow{d}\;
        N\!\left(0,\;
          \frac{\mathrm{Var}(Y_i - \beta T_i \mid Z_i = z)\,p_z^{-1}}
               {\Delta_T^2}
        \right)
      \]
      summed appropriately over $z \in \{0,1\}$, where
      $p_z = \Pr(Z = z)$.  Confirm that this matches the IV
      variance formula~\eqref{w12:eq:iv-var} in the scalar
      no-covariate case $p = 0$.
  \end{enumerate}

% ------------------------------------------------------------------
\item \textbf{Matrix form of 2SLS and why the second-stage standard
  errors are wrong.}
  Stack observations into matrices: let $\mathbf{Y} \in \mathbb{R}^n$,
  $\mathbf{D} \in \mathbb{R}^{n \times k}$ be the full regressor
  matrix (including the intercept, $T$, and $X$), and
  $\mathbf{W} \in \mathbb{R}^{n \times m}$ be the full instrument
  matrix (including the intercept, $X$, and $Z$).  Let
  $P_{\mathbf{W}} = \mathbf{W}(\mathbf{W}^\top\mathbf{W})^{-1}\mathbf{W}^\top$
  be the orthogonal projection onto the column space of $\mathbf{W}$.
  \begin{enumerate}
    \item Show that the 2SLS estimator can be written as
      \[
        \hat\theta_{\mathrm{2SLS}}
        = \bigl(\mathbf{D}^\top P_{\mathbf{W}}\mathbf{D}\bigr)^{-1}
          \mathbf{D}^\top P_{\mathbf{W}}\mathbf{Y}.
      \]
      (\emph{Hint:} the first stage replaces $\mathbf{D}$ by its
      projection $\hat{\mathbf{D}} = P_{\mathbf{W}}\mathbf{D}$; then
      apply OLS to the second-stage regression of $\mathbf{Y}$ on
      $\hat{\mathbf{D}}$, using idempotency of $P_{\mathbf{W}}$.)
    \item In the single-instrument, no-covariate case ($k = m = 2$,
      intercept plus $T$ or $Z$), verify from the matrix formula
      above that $\hat\beta_{\mathrm{2SLS}} = \hat\beta_{\mathrm{IV}}$,
      reproducing Theorem~\ref{w12:thm:equiv} without the
      Frisch--Waugh--Lovell argument.
    \item The second-stage OLS regression uses the fitted regressor
      matrix $\hat{\mathbf{D}}$ in place of $\mathbf{D}$.  Let
      $\hat{\boldsymbol{\varepsilon}}_{\mathrm{2nd}}
      = \mathbf{Y} - \hat{\mathbf{D}}\hat\theta_{\mathrm{2SLS}}$
      denote the second-stage residuals actually used by OLS.  Show
      that $\hat{\boldsymbol{\varepsilon}}_{\mathrm{2nd}}
      \neq \mathbf{Y} - \mathbf{D}\hat\theta_{\mathrm{2SLS}}$
      in general.  Explain why this discrepancy makes the
      second-stage OLS standard errors invalid, and identify the
      correct residuals that should enter the sandwich variance
      estimator~\eqref{w12:eq:sandwich-est}.
  \end{enumerate}

% ------------------------------------------------------------------
\item \textbf{Efficiency of GMM and the Sargan--Hansen $J$-statistic.}
  \begin{enumerate}
    \item Prove that $V_{\mathrm{GMM}}(\Omega) \succeq V_{\mathrm{eff}}$
      for every positive-definite $\Omega$, where
      $V_{\mathrm{eff}} = (A^\top\Sigma^{-1}A)^{-1}$ is the efficient
      GMM variance~\eqref{w12:eq:eff-var}.  (\emph{Hint:} factor
      $V_{\mathrm{GMM}}(\Omega) - V_{\mathrm{eff}}$ as a product of the
      form $C^\top \Sigma^{-1} C$ for a suitable matrix $C$ by
      completing the square; recall that $B^\top\Sigma^{-1}B \succeq 0$
      for any matrix $B$ and positive-definite $\Sigma$.)
    \item Under homoskedasticity, show that 2SLS is the efficient
      GMM estimator by verifying that the 2SLS weighting matrix
      $\Omega_{\mathrm{2SLS}} = \E[W_i W_i^\top]^{-1}$ is a scalar
      multiple of $\Sigma^{-1}$.  Conclude that no efficiency gain
      over 2SLS is possible in the homoskedastic linear IV model.
    \item Return to Example~\ref{w12:ex:gmm-toy} with two instruments
      and no covariates.  The GMM estimator under identity weighting
      gives $\hat\beta \approx 0.764$ and residual sample moments
      $\hat{U}_1 = 0.018$, $\hat{U}_2 = -0.011$.  Treating
      $\hat\Sigma = I_2$ and $n = 200$, compute the Sargan--Hansen
      $J$-statistic $J = n\,\hat{U}(\hat\beta)^\top\hat\Sigma^{-1}
      \hat{U}(\hat\beta)$ and determine, using the $\chi^2_1$ critical
      value at the 5\% level, whether the overidentifying restriction
      is rejected.  Interpret the result in terms of instrument
      validity.
  \end{enumerate}

% ------------------------------------------------------------------
\item \textbf{GEL first-order conditions and the minimum-discrepancy
  dual.}
  \begin{enumerate}
    \item For empirical likelihood (EL), $G(v) = \log(1 - v)$.
      At fixed $\theta$, write the first-order condition
      $\partial/\partial\lambda[\,n^{-1}\sum_i G(\lambda^\top
      U_i(\theta))] = 0$ for the inner supremum in~\eqref{w12:eq:gel}.
      Show that it implies $\sum_{i=1}^n \hat\pi_i U_i(\theta) = 0$,
      where $\hat\pi_i \propto (1 - \hat\lambda^\top U_i(\theta))^{-1}$.
      Verify that the normalization in~\eqref{w12:eq:gel-probs} is
      consistent with $g(v) = (1-v)^{-1}$.
    \item For the continuous updating estimator (CUE), $G(v) =
      v + v^2/2$.  Solve the inner supremum over $\lambda$
      explicitly at fixed $\theta$: show that
      \[
        \hat\lambda
        = -\Bigl[n^{-1}{\textstyle\sum_i}
          U_i(\theta)U_i(\theta)^\top\Bigr]^{-1}\bar{U}(\theta),
      \]
      and substitute back to confirm that the profile objective
      equals $\tfrac{1}{2}\bar{U}(\theta)^\top
      [n^{-1}\sum_i U_i(\theta)U_i(\theta)^\top]^{-1}\bar{U}(\theta)$,
      establishing Theorem~\ref{w12:thm:cue-gel}.
    \item In the exactly identified case ($m = k$), show that
      $\hat\lambda = 0$ at any GEL solution $\hat\theta$.  Conclude
      that every GEL estimator coincides with the just-identified GMM
      estimator in this case, and that the empirical
      probabilities~\eqref{w12:eq:gel-probs} all equal $1/n$.
      (\emph{Hint:} in the exactly identified case
      $\bar{U}(\hat\theta) = 0$ exactly; use this to argue that
      $\lambda = 0$ is the unique optimizer in the inner supremum,
      and verify that the corresponding primal weights from the
      conjugate relationship~\eqref{w12:eq:gel-conjugate} all equal
      the reference value $\hat\omega_i = g(0) = 1$.)
  \end{enumerate}

% ------------------------------------------------------------------
\item \textbf{Control function, endogeneity testing, and the limits of
  instrument diagnostics.}
  Consider the structural model $Y = \alpha + \beta T + \varepsilon$
  and first stage $T = \pi Z + \eta$, with
  $\mathrm{Cov}(\varepsilon, \eta) \neq 0$.
  Write $\varepsilon = \rho\eta + \xi$ with
  $\rho = \mathrm{Cov}(\varepsilon,\eta)/\mathrm{Var}(\eta)$.
  \begin{enumerate}
    \item Show that $\E[\xi \mid \eta] = 0$ by construction, and
      then show that $\E[\xi \mid Z] = 0$ follows from the IV
      exogeneity assumption $\E[\varepsilon \mid Z] = 0$ together
      with $\E[\eta \mid Z] = 0$.  Conclude that conditioning on
      $\eta$---or on any sufficient statistic for it---renders $T$
      exogenous in the augmented regression of $Y$ on $(T, \eta)$.
    \item Show that the coefficient on $\hat\eta$ in the augmented
      OLS regression of $Y$ on $(T, \hat\eta)$ is a consistent
      estimator of $\rho$.  Under $H_0\colon\rho = 0$ (exogeneity
      of $T$), argue that the $t$-statistic on $\hat\eta$ provides a
      valid asymptotic test of endogeneity, and connect this to the
      Hausman (1978) test.
    \item Suppose an instrument $Z$ affects wages both through
      education (the intended channel $Z \to T \to Y$) and through
      a direct network effect ($Z \to Y$), but the model is exactly
      identified ($q = 1$).  Explain why neither the first-stage
      $F$-test nor the control function endogeneity test can detect
      this exclusion-restriction violation.  What additional
      information or design feature would be needed to test for
      it?
  \end{enumerate}

\end{enumerate}

% Bibliography (standalone only)
\ifSubfilesClassLoaded{%
  \bibliographystyle{plainnat}
  \bibliography{causal_inference}
}{}

\end{document}